\documentclass[11pt]{amsart}

\usepackage{amsthm, amsfonts, amssymb, amscd, rotating}
\usepackage[pagebackref,colorlinks]{hyperref}
%\usepackage[notref,notcite]{showkeys}
\usepackage{tikz-cd}
\usepackage{geometry}
\usepackage{marginnote}
\usepackage{mathtools}

\theoremstyle{definition}
\newtheorem{ntn}{Notation}[section]
\newtheorem{dfn}[ntn]{Definition}
\theoremstyle{plain}
\newtheorem{lem}[ntn]{Lemma}
\newtheorem{prp}[ntn]{Proposition}
\newtheorem{thm}[ntn]{Theorem}
\newtheorem{cor}[ntn]{Corollary}
\newtheorem{conj}[ntn]{Conjecture}
\theoremstyle{definition}
\newtheorem{question}[ntn]{Question}
\newtheorem{rem}[ntn]{Remark}
\newtheorem{exm}[ntn]{Example}
%\newtheorem{sub}[ntn]{Subsection}

\numberwithin{equation}{section}

\newcommand{\N}{\mathbb{N}}
\newcommand{\z}{\mathbb{Z}}
\newcommand{\q}{\mathbb{Q}}
\newcommand{\R}{\mathbb{R}}
\newcommand{\C}{\mathbb{C}}
\newcommand{\F}{\mathbb{F}}

\newcommand{\EE}{\mathcal{E}}
\newcommand{\BB}{\mathcal{B}}
\newcommand{\GG}{\mathcal{G}}
\newcommand{\RR}{\mathcal{R}}
\newcommand{\WW}{\mathcal{W}}
\newcommand{\II}{\mathcal{I}}
\newcommand{\PP}{\mathcal{P}}
\newcommand{\RP}{\mathcal{RP}}
\newcommand{\RB}{\mathcal{RB}}
\renewcommand{\SS}{\mathcal{S}}

\newcommand{\ppp}{\mathfrak{p}}
\newcommand{\mmm}{\mathfrak{m}}

\renewcommand{\aa}{{A^\times}}
\newcommand{\tors}{{{\rm Tor}_1^{\z}}}
\newcommand{\tor}{{{\rm Tor}}}
\newcommand{\half}{{\left[\frac{1}{2}\right]}}
\newcommand{\third}{{\left[\frac{1}{3}\right]}}
\newcommand{\pth}{{\left[ \frac{1}{p} \right]}}
\newcommand{\mth}[1]{\left[ \frac{1}{#1} \right]}

\newcommand{\mt}{\mapsto}
\newcommand{\lan}{\langle}
\newcommand{\ran}{\rangle}
\newcommand{\se}{\subseteq}
\newcommand{\arr}{\rightarrow}
\newcommand{\larr}{\longrightarrow}
\newcommand{\harr}{\hookrightarrow}
\newcommand{\two}{\twoheadrightarrow}
\newcommand{\Lan}{\langle \! \langle}
\newcommand{\Ran}{\rangle \! \rangle}

\newcommand{\GE}{{\rm GE}}
\newcommand{\GW}{{\rm GW}}
\newcommand{\MW}{{\rm MW}}
\newcommand{\Eb}{\mathbb{E}}
\newcommand{\Ee}{\mathrm{E}}
\newcommand{\Ind}{{\rm Ind}}
\newcommand{\Hom}{{\rm Hom}}
\newcommand{\Tot}{{\rm Tot}}
\newcommand{\stabe}{{\rm Stab}}
\newcommand{\GL}{\mathit{{\rm GL}}}
\newcommand{\SL}{\mathit{{\rm SL}}}
\newcommand{\GM}{\mathit{{\rm GM}}}
\newcommand{\SM}{\mathit{{\rm SM}}}
\newcommand{\PGL}{\mathit{{\rm PGL}}}
\newcommand{\PGE}{\mathit{{\rm PGE}}}

\newcommand{\id}{{\rm id}}
\newcommand{\im}{{\rm im}}
\newcommand{\inc}{{\rm inc}}
\newcommand{\ind}{{\rm ind}}
\newcommand{\sign}{{\rm sign}}
\newcommand{\diag}{{\rm diag}}
\newcommand{\rank}{{\rm rank}}
\renewcommand{\char}{{\rm char}}
\newcommand{\coker}{{\rm coker}}

\newcommand {\mtx}[2]
{\left(\!\!\!
	\begin{array}{cc}
		#1 \\
		#2 
	\end{array}
	\!\!\!\right)}

\newcommand{\mtxx}[4]
{\left(\!\!
	\begin{array}{cc}
		\!\!#1 & \!\!#2\\
		\!\!#3 & \!\!#4
	\end{array}\!\!
	\right)}


\begin{document}
	\title{Remarks on refined scissors congruence group and the third homology of \texorpdfstring{$\SL_2$}{Lg}}
	\author{Elvis Torres P\'erez}
	\address{\sf 
		Department of Sciences, Pontifical Catholic University of Peru (PUCP), 
		Lima, Peru}
	\email{etorresp@pucp.edu.pe}
	
	\begin{abstract}
		This work revisits the key sufficient conditions on a ring $A$ that guarantee the existence of the exact sequence
		\[
		H_3(\SM_2(A),\z)\to H_3(\SL_2(A),\z)\to\RB(A)\to 0
		\]
		and of the isomorphism 
		\[
		H_3(\SL_2(A),\SM_2(A);\z)\simeq\RP_1(A).
		\]
		We show that there are rings which satisfy the relations above but do not satisfy the condition $-1\in\aa^2$, for example the rings $\z\half$ and $\z\mth{10}$. In addition, we study the special case of non-archimedean local fields, obtaining a refined Bloch-Wigner exact sequence when $-1\notin\aa^2$.
	\end{abstract}
	\maketitle
	%%%%%%%%%%%%%%%%%%%%%%%%%%%%%%%%%%%%%%%%%%%%%%%%%%%%%
	\section{Introduction}
	%%%%%%%%%%%%%%%%%%%%%%%%%%%%%%%%%%%%%%%%%%%%%%%%%%%%%
	%The Bloch group $\BB(\C)$ was defined by S. Bloch in {\bf refff} to approximates the $K_3^\ind(\C)$, this group (in Suslin words) describes the non-trivial relations between tensors $x\otimes(1-x)$ where $x\neq 0,1$. 
	In \cite{suslin1991}, Andrei A. Suslin states and proves a Bloch-Wigner exact sequence for any infinite field $F$: 
	\begin{equation}\label{sus-BW-seq}
		0\to\tors(\mu(F),\mu(F))^{\tilde{}}\to K^\ind_3(F)\to \BB(F)\to 0.    
	\end{equation}
	Here $\BB(F)$ is the so-called Bloch group, which (in Suslin's words) describes the non-trivial relations among the tensors $x\otimes(1-x)$ with $x\neq 0,1$; it was defined by S. Bloch in \cite{bloch2000} precisely as an approximation to $K_3^\ind(F)$. The group $K^\ind_3(F)$ is the indecomposable part of $K_3(F)$, and $\tors(\mu(F),\mu(F))^{\tilde{}}$ is the unique non-trivial extension of $\tors(\mu(F),\mu(F))$ by $\z/2$. As an intermediate result, Suslin obtains the exact sequence
	\begin{equation}\label{gm2-gl2-seq}
		H_3(\GM_2(F),\z)\to H_3(\GL_2(F),\z)\to \BB(F)\to 0.    
	\end{equation}
	
	Some time ago, C. Sah proved that if $F$ is algebraically closed, then $K^\ind_3(F)$ coincides with the group $H_3(\SL_2(F),\z)$. In 2013, with the aim of understanding the relation between $H_3(\SL_2(F),\z)$ and $K_3^\ind(F)$, Kevin Hutchinson \cite{hutchinson-2013} defined the refined Bloch group $\RB(F)$ as a subgroup of $\RP_1(F)$, the so-called {\it refined scissors congruence group} (the name was inspired by the work of J. L. Dupont and C. Sah on Hilbert's third problem in hyperbolic geometry \cite{dupont-sah1982}), and constructed a refined Bloch-Wigner complex (for infinite fields)
	\[
	0\to \tors(\mu(F),\mu(F))\to H_3(\SL_2(F),\z)\to\RB(F)\to 0
	\]
	which is exact at the outer terms and whose homology at the middle term is annihilated by $4$.
	This result was generalized by Hutchinson to local rings with sufficiently large residue field \cite{hutchinson2017}. In \cite{B-E--2023}, while studying the above complex, B. Mirzaii and the author obtained, for a commutative ring $A$, an $\SL_2$-version of Suslin's exact sequence \eqref{gm2-gl2-seq}: 
	\begin{equation}\label{sm2-sl2-seq}
		H_3(\SM_2(A),\z)\to H_3(\SL_2(A),\z)\to\RB(A)\to 0
	\end{equation}
	when $A$ is a universal $\GE_2$-ring (see definitions below) that satisfies
	\begin{enumerate}
		\item $\mu_2(A)=\{\pm 1\}$ and $-1\in\aa^2$.
		\item $H_i(B(A),\z)\simeq H_i(T(A),\z)$ for $i=2,3$.
	\end{enumerate}
	In \cite[Remark 6.7]{B-E--2023} we suggested that the condition $-1\in\aa^2$ is in fact not essential (at least when $A$ is a domain).
	The present paper studies in more depth some of the key conditions implying \eqref{sm2-sl2-seq}, as well as other results of \cite{B-E--2023}, such as
	\begin{equation}\label{sl2-sm2-rel}
		H_3(\SL_2(A),\SM_2(A);\z)\simeq \RP_1(A).
	\end{equation}
	
	More precisely, in \cite[Corollary 6.4]{B-E--2023} we proved that if $-1\in\aa^2$, then certain maps $\alpha$ and $\gamma$ have the same image; here we study the consequences of the partial inclusions $\im(\alpha)\subseteq\im(\gamma)$ and $\im(\gamma)\subseteq\im(\alpha)$. The first main theorem describes the relation between these two inclusions when $A$ is a domain. The second main theorem provides a generalization of the exact sequence \eqref{sm2-sl2-seq} when $\im(\gamma)\subseteq\im(\alpha)$; when the two images coincide, we recover \eqref{sm2-sl2-seq} itself.
	As corollaries, we obtain a slight generalization of Theorem 6.6 of \cite{B-E--2023}, namely the exact sequence
	\[
	\II_A\otimes\mu_2(A)\to H_3(\SL_2(A),\SM_2(A);\z)\to \frac{\RP_1(A)}{\II_A\psi_1(-1)}\to 0.
	\]
	as well as the isomorphism
	\[
	H_3(\SL_2(A),\SM_2(A);\z)\simeq \RP_1(A)
	\]
	when $-1\in\aa^2$ (Theorem 8.2 of \cite{B-E--2023}). In fact, we show that in all these results the condition $H_2(B(A),\z)\simeq H_2(T(A),\z)$ is not needed.
	
	Some straightforward examples are the rings whose group of square classes has order $2$, such as $\z$, the rings $\z/p^r$ with $p$ an odd prime, and the finite fields $\F_q$. Less immediate examples are the rings of the form $\z\mth{m}$; we show that $\z\mth{2}$ and $\z\mth{10}$ satisfy all the hypotheses of our main theorems (see Remark \ref{lack-on-zm}). Finally, the last section deals with non-archimedean, non-dyadic local fields and establishes the exact sequence \eqref{sm2-sl2-seq} when $-1\notin\aa^2$; this completes the discussion of non-dyadic, non-archimedean local fields of characteristic different from $2$. Moreover, in this case we obtain a refined Bloch-Wigner exact sequence, determining along the way the complete structure of $H_3(\SM_2(F),\z)$. Namely, we obtain the exact sequence
	\[
	0\to\mu(F)^{\tilde{}}\to H_3(\SL_2(F),\z)\to \RB(F)\to 0
	\]
	
	{\bf Notation.} In this article all rings are commutative (with the possible exception of group rings) and have a unit element $1$. For a ring $A$, $\aa$ denotes the group of invertible elements of $A$. 
	Let $\WW_A$ denote the set of all $a\in \aa$ such that $1-a\in \aa$; explicitly,
	\[
	\WW_A:=\{a\in A: a(1-a)\in \aa\}.
	\]
	Let $\GG_A:=\aa/(\aa)^2$ and $\RR_A:=\z[\GG_A]$. The element of $\GG_A$ represented by 
	$a \in \aa$ is denoted by $\lan a \ran$. We set $\Lan a\Ran:=\lan a\ran -1\in \RR_A$. Finally, let $\epsilon:\RR_A\to\z$ be the homomorphism defined by $\lan a\ran\mapsto1$ for every $a\in\aa$, and let $\II_A$ denote its kernel.\\
	~\\
	%{\bf Acknowledgements.}
	%The second author acknowledges that during work on this article he was supported by CAPES (Coordena\c{c}\~ao de
	%Aperfei\c{c}oamento de Pessoal de N\'ivel Superior) PhD fellowship (grant number 88887.357813/2019-00).
	\section{\texorpdfstring{$\GE_2$}{GE2} rings, unimodular vectors and scissors congruence}
	Let $A$ be a ring. The matrices
	\[
	E(a):=\mtxx{a}{1}{-1}{0}
	\]
	generate a subgroup of $\SL_2(A)$ which coincides with the subgroup $\Ee_2(A)$ generated by the elementary matrices
	\[
	E_{12}(a):=\mtxx{1}{a}{0}{1},\quad \quad E_{21}(a):=\mtxx{1}{0}{a}{1}, 
	\]
	as one sees from the formulas
	\[
	E_{12}(a)=E(-a)E(0)^{-1}, \quad E_{21}(a)=E(0)^{-1}E(a).
	\]
	
	The subgroup $\Ee_2(A)$, together with the subgroup $D_2$ generated by the diagonal matrices in $\GL_2(A)$, generates the subgroup $\GE_2(A)\subseteq \GL_2(A)$. It is known that if $A$ is a field or a Euclidean domain, then $\GE_2(A)=\GL_2(A)$; rings with this property are called $\GE_2$-rings. It is not difficult to show that, for any ring, $\Ee_2(A)=\SL_2(A)$ if and only if $A$ is a $\GE_2$-ring. In \cite{cohn1966}, Cohn lists the following relations among the matrices $E(a)$:
	\begin{enumerate}
		\item For any $u,v\in\aa$, $D(uv)=D(u)D(v)$, where $D(a)=E(-a)E(-a^{-1})E(-a)$ for $a\in\aa$.
		\item For any $a,b\in A$, $E(a)E(0)E(b)=-D(-1)E(a+b)$.
		\item For any $u\in\aa$ and $a\in A$, $D(u)E(a)D(u)=E(u^2a)$.
	\end{enumerate}
	
	He also observed that for many rings these relations are {\it universal}, meaning that every relation among the matrices of $\Ee_2(A)$ is a consequence of Cohn's relations; such rings are called {\it universal for $\GE_2$}. A $\GE_2$-ring which is universal for $\GE_2$ is called a {\it universal $\GE_2$-ring}; see \cite[\S 4]{hutchinson2022} for a more explicit definition of these notions.
	
	Let $A$ be a ring. A (column) vector $\pmb{u}=\mtx{u_1}{u_2}\in A^2$ is said to be \textit{unimodular} if there exists a vector $\pmb{v}=\mtx{v_1}{v_2}$ such that the matrix $(\pmb{u},\pmb{v}):=\mtxx{u_1}{v_1}{u_2}{v_2}$ lies in $\GL_2(A)$. Consider the free abelian group $X_k(A^2)$ generated by the natural $\GL_2(A)$-set (with the diagonal action)
	\[
	\{(\lan\pmb{v_0}\ran,\lan\pmb{v_1}\ran,\cdots,\lan\pmb{v_k}\ran): \pmb{v_i} \text{ is unimodular and } (\pmb{v_i},\pmb{v_j})\in \GL_2(A) \text{ for } i\neq j\}
	\]
	where $\lan \pmb{v}\ran$ denotes the line generated by the vector $\pmb{v}$. Among these lines we single out $\pmb{\infty}=\left\lan\mtx{1}{0}\right\ran$ and $\pmb{0}=\left\lan\mtx{0}{1}\right\ran$, as well as the line $\pmb{a}=\left\lan\mtx{1}{a}\right\ran$ for $a\neq 0$. In this way $X_k(A^2)$ is a left $\GL_2(A)$-module (and hence an $\SL_2(A)$-module), which can be turned into a right module by setting $m\cdot g:=g^{-1}\cdot m$ for $m\in X_k(A^2)$ and $g\in\GL_2(A)$.
	
	For $k\geq 1$, consider the $k$-th differential
	\[
	\partial_k:X_k(A^2)\to X_{k-1}(A^2)
	\]
	defined by
	\[
	\partial_k((\lan\pmb{v_0}\ran,\lan\pmb{v_1}\ran,\cdots,\lan\pmb{v_k}\ran))=\sum^k_{i=0}(-1)^i(\lan\pmb{v_0}\ran,\lan\pmb{v_1}\ran,\cdots,\hat{\lan\pmb{v_i}\ran},\cdots,\lan\pmb{v_k}\ran)
	\]
	where $\hat{\lan\pmb{v_i}\ran}$ indicates that the $i$-th component is omitted. Let $\partial_0:X_0(A^2)\to\z$ be defined by $\lan\pmb{v_0}\ran\mapsto 1$. We thus obtain the complex
	\[
	\begin{tikzcd}
		X_{\bullet}(A^2)\to\z:& \cdots \ar[r] & X_2(A^2)\ar[r,"\partial_2"] & X_1(A^2)\ar[r,"\partial_1"] & X_0(A^2) \ar[r,"\partial_0"] & \z \ar[r] & 0.
	\end{tikzcd}
	\]
	We shall say that the above complex is \textit{exact in dimension} $<k$ if the complex
	\[
	\begin{tikzcd}
		X_k(A^2)\ar[r,"\partial_k"] & X_{k-1}(A^2)\ar[r,"\partial_{k-1}"] & \cdots \ar[r,"\partial_2"] & X_1(A^2)\ar[r,"\partial_1"] & X_0(A^2) \ar[r,"\partial_0"] & \z \ar[r] & 0
	\end{tikzcd}
	\]
	is exact. For instance, for local rings (or fields), the exactness of $X_{\bullet}(A^2)\to\z$ is governed by the number of elements of the residue field.
	
	%\begin{prp}[Hutchinson]
	%    Let $A$ a local ring. If $k$ is the number of elements of the residue field, then the complex $X_{\bullet}(A^2)\to\z$ is exact in dimension $<k$.
	%\end{prp}
	%\begin{proof}
	%    See \cite[Theorem 3.21]{hutchinson2017}.
	%\end{proof}
	In \cite{hutchinson2022}, Hutchinson studies the complex of unimodular vectors as the clique complex of a graph $\Gamma(A)$ whose vertices are the unimodular rows and in which an edge joins $\lan\pmb{u}\ran$ and $\lan\pmb{v}\ran$ whenever $[\pmb{u},\pmb{v}]\in \GL_2(A)$. In this way he obtains a description of the homology groups $H_i(X_\bullet(A^2))$ for $i=0,1$. 
	\begin{thm}[Hutchinson]
		For any commutative ring $A$, we have the isomorphism
		\[
		H_0(X_\bullet(A^2))\simeq\z[\SL_2(A)/\Ee_2(A)]
		\]
	\end{thm}
	\begin{proof}
		See \cite[Theorem 2.3]{B-E--2026}.
	\end{proof}
	Let $K_2(2,A)$ be the rank one $K_2$-group and let $C(2,A)$ be the subgroup of $K_2(2,A)$ generated by certain symbols $c(u,v)$. It is known that $A$ is universal for $\GE_2$ if and only if $K_2(2,A)$ is generated by these symbols; see \cite[Appendix A]{hutchinson2022} for a complete proof. In this setting, Hutchinson describes $H_1(X_\bullet(A^2))$ as follows.
	\begin{thm}[Hutchinson]
		For any commutative ring $A$, we have the isomorphism
		\[
		H_1(X_\bullet(A^2))\simeq \Ind^{\SL_2(A)}_{\Ee_2(A)}\left(\frac{K_2(2,A)}{C(2,A)}\right)^{\textrm{ab}}
		\]
	\end{thm}
	\begin{proof}
		See \cite[Theorem 2.5]{B-E--2026}.
	\end{proof}
	The theorems above imply that if $A$ is a universal $\GE_2$-ring, then $X_\bullet(A^2)\to\z$ is exact in dimension $<2$.
	
	\begin{exm}
		%\label{ex-ge2-rings}
		\begin{enumerate}
			\item \cite[\S 4]{cohn1966} Local rings (fields) are universal $\GE_2$-rings.
			\item \cite[Example 4.7]{C-H2022} $\z$ is a Euclidean domain and $K_2(2,\z)$ is generated by $c(-1,-1)$; hence $\z$ is a universal $\GE_2$-ring.
			\item Let $d$ be a squarefree positive integer and let $\mathcal{O}_d$ be the ring of integers of $\q[\sqrt{-d}]$. Then $\mathcal{O}_d$ is not a $\GE_2$-ring unless $d=1,2,3,7,11$ \cite[Theorem 6.1]{cohn1966}, and only in the cases $d=1,3$ is it universal for $\GE_2$ \cite[Remarks]{cohn1968}.
			\item \cite[Example 4.8]{C-H2022} The Euclidean domains $\z\left[\frac{1}{m}\right]$, where $m=p_1^{\alpha_1}\cdots p_r^{\alpha_r}$ and $(\z/p_i)^\times$ is generated by $\{-1,p_1,\dots,p_{i-1}\}$ for all $i\leq r$. In particular, $\z\half$, $\z\left[\frac{1}{10}\right]$ and $\z\left[\frac{1}{290}\right]$ are universal $\GE_2$-rings.
		\end{enumerate}
	\end{exm}
	
	For any $k\geq 0$, consider the subgroup $Z_k(A^2):=\ker(\partial_k)$ (we shall sometimes write $X_k$ and $Z_k$ instead of $X_k(A^2)$ and $Z_k(A^2)$). Following Coronado and Hutchinson \cite{C-H2022}, we define the group $\RP(A):=H_0(\SL_2(A),Z_2(A^2))=Z_2(A^2)_{\SL_2(A)}$. The natural action of $\GL_2(A)$ on this group (by conjugation) endows $\RP(A)$ with the structure of an $\RR_A$-module (via the matrix $\mtxx{1}{0}{0}{a}$ attached to $a\in\aa$). The inclusion map $Z_2(A^2)\to X_2(A^2)$ induces a map $\lambda_1:\RP(A)\to X_2(A^2)_{\SL_2(A)}$, and since $X_2(A^2)=\bigoplus_{\lan a\ran\in\GG_A}\Ind_{\mu_2(A)}^{\SL_2(A)}\z$, Shapiro's lemma gives $X_2(A^2)_{\SL_2(A)}\simeq\RR_A$. The \textit{refined scissors congruence group} is then defined as the group $\RP_1(A):=\ker(\lambda_1:\RP(A)\to\RR_A)$. 
	
	For $a\in\aa$, consider the element $\psi_1(a)\in\RP(A)$ represented by
	\[
	(\pmb{\infty},\pmb{0},\pmb{a})+(\pmb{0},\pmb{\infty},\pmb{a})-(\pmb{\infty},\pmb{0},\pmb{1})-(\pmb{0},\pmb{\infty},\pmb{1})
	\]
	Coronado and Hutchinson establish the following properties of this element in \cite{C-H2022}.
	\begin{lem}\label{lem-psi1}
		Let $A$ be a ring. The map $\psi_1:\aa\to\RP(A)$ satisfies the following properties:
		\begin{enumerate}
			\item $\psi_1(ab)=\lan a\ran\psi_1(b)+\psi_1(a)$.
			\item $\lan -1\ran\psi_1(a)=\psi_1(a^{-1})$.
			\item $\Lan a\Ran\psi_1(b)=\Lan b\Ran\psi_1(a)$.
			\item $\psi_1(ab^2)=\psi_1(a)+\psi_1(b^2)$.
			\item $2\psi_1(-1)=0$.
			\item $\psi_1(a^2)=\Lan a\Ran\psi_1(-1)=\Lan -1\Ran\psi_1(a)$.
			\item $2\psi_1(a^2)=0$ for all $a\in\aa$. If $-1\in\aa^2$, then $\psi_1(a^2)=0$ for all $a\in\aa$.
		\end{enumerate}
	\end{lem}
	\begin{proof}
		See \cite[Lemma 7.8 and Lemma 7.25]{C-H2022}.
	\end{proof}
	Consider the truncated complex $X^{\tau}_{\bullet}(A^2)\to\z$ given by $X^{\tau}_{k}(A^2)=X_{k}(A^2)$ for $k<3$, $X^{\tau}_{3}(A^2)=Z_2(A^2)$ and $X^{\tau}_{k}(A^2)=0$ for $k>3$. Tensoring this complex with a projective resolution $F_{\bullet}\to\z$ over $\RR_A$, we obtain the double complex $F_{\bullet}\otimes_{\SL_2(A)}X^{\tau}_{\bullet}$ and a spectral sequence
	\[
	E^1_{p,q}(\SL_2(A),X^{\tau})=\begin{cases}
		H_q(\SL_2(A),X_p(A^2)) & p<3\\
		H_q(\SL_2(A),Z_2(A^2)) & p=3\\
		0 & p>3
	\end{cases}\Longrightarrow H_{p+q}(\SL_2(A),X^{\tau}_{\bullet}(A^2)).
	\]
	The first page of this spectral sequence is described by Coronado and Hutchinson in \cite[Lemma 3.2]{C-H2022}; its differentials, which we denote by $d^r_{p,q}(X^\tau)$, are known:
	\begin{equation}
		\begin{tikzcd}
			H_3(B(A),\z) & \ar[l,"(w - \id)_*"'] H_3(T(A),\z) & * & * \\
			\aa\wedge\aa\oplus\SS_2 & \ar[l,"0"'] \aa\wedge\aa & \ar[l] 0 & \ar[l] H_2(\SL_2(A,Z_2) \\
			H_1(B(A),\z) & \ar[l,"(\cdot)^{-2}"'] \aa & \ar[l,"\epsilon\otimes\inc"'] \RR\otimes\mu_2(A) & \ar[l,"\inc_*"'] H_1(\SL_2(A),Z_2)\\
			\z &\ar[l,"0"'] \z & \ar[l,"\epsilon"'] \RR & \ar[l,"\inc_*=\lambda_1"'] \RP(A) 
		\end{tikzcd}
	\end{equation}
	We see that $\RP(A)=E^1_{3,0}(\SL_2(A),X^{\tau})$, and since $E^2_{1,1}(\SL_2(A),X^\tau)=0$, we have $\RP_1(A)=E^3_{3,0}(\SL_2(A),X^{\tau})$. The \textit{refined Bloch group} of $A$ is defined as $\RB(A):=\ker(\bar{d}^3_{3,0}:\RP_1(A)\to E^3_{0,2}(\SL_2(A),X^{\tau}))$. Moreover, if the complex $X^{\tau}(A^2)\to\z$ is exact, then the spectral sequence above converges to $H_{p+q}(\SL_2(A),\z)$. See \cite[\S 3]{C-H2022} for further details.
	%%%%%%%%%%%%%%%%%%%%%%%%%%%%%%%%%%%%%%%%%%%%%%%%%%%%%%%%%%%%%%%%%%%%
	\section{A spectral sequence for \texorpdfstring{$\SM_2$}{SM2} and \texorpdfstring{$\SL_2$}{SL2}}
	%%%%%%%%%%%%%%%%%%%%%%%%%%%%%%%%%%%%%%%%%%%%%%%%%%%%%%%%%%%%%%%%%%%
	Let $A$ be a ring and consider the following map of complexes.
	\begin{equation}
		\label{sm-sl-cplx}
		\begin{tikzcd}
			0\ar[r] & \hat{Z}_1(A^2)\ar[r]\ar[d,"\inc"] & \hat{X}_1(A^2)\ar[r,"\hat{\partial}_1"]\ar[d,"\inc"] & \hat{X}_0(A^2)\ar[r,"\hat{\partial}_0"]\ar[d,"\inc"] & \z\ar[d,equal]\\
			0\ar[r] & Z_1(A^2)\ar[r] & X_1(A^2)\ar[r,"\partial_1"] & X_0(A^2)\ar[r,"\partial_0"] & \z
		\end{tikzcd}
	\end{equation}
	where the groups in the upper row $\hat{X}_0(A^2)=\z\{(\pmb{\infty}),(\pmb{0})\}$, $\hat{X}_1(A^2)=\z\{(\pmb{\infty},\pmb{0}),(\pmb{0},\pmb{\infty})\}$ and $\hat{Z}_1(A^2)=\z\{(\pmb{\infty},\pmb{0})+(\pmb{0},\pmb{\infty})\}\simeq\z$ are modules over the group of monomial matrices
	\[
	\SM_2(A):=\left\{\mtxx{a}{0}{0}{a^{-1}},\mtxx{0}{b^{-1}}{-b}{0}: a,b\in\aa\right\}.
	\]
	and the lower row is a truncated piece of the complex $X_{\bullet}(A^2)\to\z$. It is not difficult to show that
	\begin{equation}
		\label{struct-X-hatX}
		\hat{X}_0\simeq\Ind^{\SM_2(A)}_{T(A)}\z,\quad \hat{X}_1\simeq\Ind^{\SM_2(A)}_{T(A)}\z,\quad  X_0\simeq\Ind^{\SM_2(A)}_{B(A)}\z,\quad X_1\simeq\Ind^{\SM_2(A)}_{T(A)}\z    
	\end{equation}
	where
	\[
	T(A)=\left\{\mtxx{a}{0}{0}{a^{-1}}: a\in\aa\right\},\quad B(A)=\left\{\mtxx{a}{b}{0}{a^{-1}}: a\in\aa, b\in A\right\},
	\]
	and by Shapiro's lemma we obtain $\hat{E}^1_{0,q}\simeq H_q(T(A),\z)$, $\hat{E}^1_{1,q}\simeq H_q(T(A),\z)$ and $\hat{E}^1_{2,q}\simeq H_q(\SM_2(A),\z)$.
	
	Now, if $X_{\bullet}(A^2)\to\z$ is exact in dimension $<1$ (for example, when $A$ is a $\GE_2$-ring), tensoring \eqref{sm-sl-cplx} with a projective resolution yields the following map of spectral sequences.
	\begin{equation}\label{map-s-seq}
		\begin{tikzcd}
			\hat{E}^1_{p,q}\ar[r,Rightarrow]\ar[d,"\inc_*"] & H_{p+q}(\SM_2(A),\z)\ar[d,"\inc_*"]\\
			E^1_{p,q}\ar[r,Rightarrow] & H_{p+q}(\SL_2(A),\z).
		\end{tikzcd}
	\end{equation}
	
	Both spectral sequences were studied in \cite[\S 1 and \S 5]{B-E--2023}; note that on the $E^2$-page we have $\hat{E}^2_{2,1}\simeq\GG_A$. 
	
	Furthermore, if the complex $X_{\bullet}(A^2)\to\z$ is exact in dimension $<2$ (for example, when $A$ is a universal $\GE_2$-ring), we obtain the short exact sequence $0\to Z_2(A^2)\to X_2(A^2) \overset{\partial_2}{\to} Z_1(A^2)\to 0$, which induces the long exact sequence
	\begin{equation}\label{long-eseq}
		H_2(\SL_2(A),Z_1) \to H_1(\SL_2(A),Z_2) \to H_1(\SL_2(A),X_2)\to H_1(\SL_2(A),Z_1)\overset{\delta}{\to} \cdots.
	\end{equation}
	
	Using the structure of $X_2(A^2)$ described in \eqref{struct-X-hatX}, together with Shapiro's lemma, we obtain $H_1(\SL_2(A),X_2(A^2))\simeq\RR_A\otimes\mu_2(A)$. Truncating the exact sequence above, we get
	\[
	\begin{tikzcd}
		H_1(\SL_2(A),Z_2) \ar[r] & \RR_A\otimes\mu_2(A)\ar[r,"\gamma"] & E^1_{2,1} \ar[r,"\delta"] & \RP_1(A)\ar[r] & 0.
	\end{tikzcd}
	\]
	where $\gamma$ is induced by $\partial_2$. As in \cite[\S 4]{B-E--2023}, the maps $d^1_{2,1}:E^1_{2,1}\to\mu_2(A)\subset E^1_{1,1}$ and $d^1_{2,1}\circ\gamma=\epsilon\otimes\id:\RR_A\otimes\mu_2(A)\to\mu_2(A)$ fit into a commutative diagram, and the snake lemma then yields
	\[
	\begin{tikzcd}
		&\II_A\otimes\mu_2(A)\ar[r,"\gamma"] & E^2_{2,1} \ar[r,"\delta"] & \RP_1(A)\ar[r] & 0 
	\end{tikzcd}
	\]
	Finally, the map of spectral sequences yields the following diagram.
	\begin{equation}\label{dia-gam-alpha}
		\begin{tikzcd}
			& \hat{E}^2_{2,1}\simeq\GG_A\ar[d,"\alpha:=\inc_*"] & &\\
			\II_A\otimes\mu_2(A)\ar[r,"\gamma"] & E^2_{2,1} \ar[r,"\delta"] & \RP_1(A)\ar[r] & 0.
		\end{tikzcd}
	\end{equation}
	%%%%%%%%%%%%%%%%%%%%%%%%%%%%%%%%%%%%%%%%%%%%%%%%%%%%
	\section{The element \texorpdfstring{$\chi_1$}{X1}}
	%%%%%%%%%%%%%%%%%%%%%%%%%%%%%%%%%%%%%%%%%%%%%%%%%%%%
	From now on, the elements $(\pmb{\infty},\pmb{0},\pmb{a})$ and $(\pmb{0},\pmb{\infty},\pmb{a})$ in $X_2(A^2)$, and $(\pmb{\infty},\pmb{0})+(\pmb{0},\pmb{\infty})$ in $X_1(A^2)$, will be denoted by $X_a$, $X'_a$ and $Y$, respectively; in bar notation, the matrices $\mtxx{0}{1}{-1}{0}\in\SM_2(A)$ and $\mtxx{a}{0}{0}{a^{-1}}$, for $a\in\aa$, will be denoted by $w$ and $a$, respectively. Finally, the maps $\alpha$, $\gamma$ and $\delta$ are those of diagram \eqref{dia-gam-alpha}, and $d^2_{2,1}:E^2_{2,1}\to H_2(B(A),\z)\simeq H_2(T(A),\z)\oplus\SS_2$ (see Section \ref{sec-hb-ht} for this decomposition) is the $E^2$-differential of the spectral sequence $E^1_{p,q}$ appearing in diagram \eqref{map-s-seq}. 
	
	Consider the element $\chi_1\in E^2_{2,1}$ represented by the cycle $[w]\otimes\partial_2(X_{-1}-X_1)$.
	\begin{lem}\label{alpha-gam-chi}
		The element $\chi_1\in E^2_{2,1}$ and the maps $\alpha$, $\gamma$ and $\delta$ satisfy the following properties:
		\begin{enumerate}
			\item $\delta(\chi_1)=\psi_1(-1)$.
			\item $\chi_1=[w]\otimes Y-\gamma(\lan 1\ran\otimes (-1))$.
			\item $\gamma(\Lan a\Ran\otimes (-1))-\alpha(\lan a\ran)=\Lan a\Ran\chi_1$.
			\item $(\delta\circ\alpha)(\lan a\ran)=\psi_1(a^2)=\Lan a\Ran\psi_1(-1)$.
			\item $d^2_{2,1}\circ\alpha(\lan a\ran)=(-1\wedge a,0)$.
			\item $\gamma(\Lan -1\ran\otimes(-1))=\alpha(\lan -1\ran)$.
			\item $\Lan -1\Ran\chi_1=0$.
		\end{enumerate} 
	\end{lem}
	\begin{proof}
		\begin{enumerate}
			\item Take a preimage of $[w]\otimes \partial_2(X_{-1}-X_1)$ under $\id\otimes\partial_2$, namely $[w]\otimes(X_{-1}-X_1)$. Applying $d_1\otimes\id$, we obtain (recall that $\psi_1(-1)$ is $2$-torsion by Lemma \ref{lem-psi1})
			\begin{align*}
				(w[]-[])\otimes(X_{-1}-X_1)&=[]\otimes(X'_1-X'_{-1}-X_{-1}+X_1)\\
				&=-\psi_1(-1)\\
				&=\psi_1(-1).
			\end{align*}
			\item 
			\begin{align*}
				[w]\otimes\partial_2(X_{-1}-X_1)&=[w]\otimes\partial_2(X_{-1}+X'_{-1}-(X'_{-1}+X_1))\\
				&=[w]\otimes Y-[w]\otimes\partial_2(wX_1+X_1)\\
				&=[w]\otimes Y-(w[w]+[w])\otimes\partial_2(X_1)\\
				&=[w]\otimes Y-[-1]\otimes\partial_2(X_1)\\
				&=[w]\otimes Y-\gamma(\lan 1\ran\otimes(-1)).
			\end{align*}
			\item It is clear that $\Lan a\Ran\chi_1$ is represented by $\Lan a\Ran([w]\otimes(X_{-1}-X_1))$; the identity above then implies
			\begin{align*}
				\Lan a\Ran([w]\otimes Y-[-1]\otimes\partial_2(X_1))&=[wa]\otimes Y - [w]\otimes Y-\gamma(\Lan a\Ran\otimes (-1))\\
				&=[a]\otimes Y-\gamma(\Lan a\Ran\otimes (-1))\\
				&=\alpha(\lan a\ran)-\gamma(\Lan a\Ran\otimes (-1)).
			\end{align*}
			\item Using item (1), we have
			\begin{align*}
				\Lan a\Ran\psi_1(-1)&=\Lan a\Ran\delta(\chi_1)\\
				&=\delta(\Lan a\Ran\chi_1)\\
				&=\delta(\alpha(\lan a\ran)-\gamma(\Lan a\Ran\otimes(-1))\\
				&=\delta(\alpha(\lan a\ran))\\
				&=\delta\circ\alpha(\lan a\ran).
			\end{align*}
			\item The element $\lan a\ran\in\GG_A$ is represented by $[a]\otimes Y\in B_2(\SL_2(A))\otimes_{\SL_2(A)}Z_1(A^2)$. We have that
			\[
			[a]\otimes Y=d_2\otimes\id(([w|a]-[a^{-1}|w]+[a^{-1}|a])\otimes(\pmb{\infty},\pmb{0}))
			\]
			Applying $\id\otimes\partial_1$, we obtain, modulo $\im(d_3\otimes\id)$,
			\[
			\begin{split}
				d^2_{2,1}(\alpha(\lan a\ran))&=\{w([w|a]-[a^{-1}|w]+[a^{-1}|a])-([w|a]-[a^{-1}|w]+[a^{-1}|a])\}\otimes(\pmb{\infty})\\
				&=([-1|a]-[a|-1])\otimes(\pmb{\infty})
			\end{split}
			\]
			where the last line is obtained by adding the trivial element
			\[
			d_3\otimes\id(\{-[w|w|a]+[w|a^{-1}|a]-[w|a^{-1}|a]-[a|w|w]+[a|w|a]-[a|a^{-1}|w]-[a^{-1}|a|a^{-1}]\}\otimes(\pmb{\infty}))
			\]
			\item The element $\gamma(\Lan -1\Ran\otimes(-1))$ is represented by $[-1]\otimes\partial_2(X_{-1}-X_1)$; hence
			\begin{align*}
				[-1]\otimes\partial_2(X_{-1}-X_1)&=[-1]\otimes\partial_2(X_{-1}+X'_{-1})-[-1]\otimes\partial_2(wX_1+X_1)\\
				&=[-1]\otimes Y - 2[-1]\otimes(X_1)\\
				&=[-1]\otimes Y
			\end{align*}
			\item This is a direct consequence of (3) and (6).
		\end{enumerate}
	\end{proof}
	\begin{lem}\label{d2-chi}
		The image $d^2_{2,1}(\chi_1)$ is represented by the cycle $(\left[\mtxx{1}{1}{0}{1}\bigg|-1\right]-\left[-1\bigg|\mtxx{1}{1}{0}{1}\right])\otimes(\pmb{\infty})$; consequently, $d^2_{2,1}(\chi_1)\in\SS_2$.
	\end{lem}
	\begin{proof}
		To this end, we pass through the diagram
		\[
		\begin{tikzcd}
			B_2(\SL_2(A))\otimes_{\SL_2(A)} X_0(A^ 2) & B_2(\SL_2(A))\otimes_{\SL_2(A)} X_1(A^2) 
			\ar["\id_{B_2}\otimes \partial_1"',  l] \ar[d, "d_2\otimes \id_{X_1}"] & \\
			& B_1(\SL_2(A))\otimes_{\SL_2(A)} X_1(A^ 2) & B_1(\SL_2(A))\otimes_{\SL_2(A)} 
			Z_1(A^2).\ar[l, "\id_{B_1}\otimes \inc"']
		\end{tikzcd}
		\]
		By Lemma \ref{alpha-gam-chi} above, $[w]\otimes Y-[-1]\otimes\partial_2(X_1)$ is a representative of $\chi_1$; passing through the inclusion, we obtain
		\begin{align*}
			[w]\otimes Y-[-1]\otimes\partial_2(X_1)=& (w[w]+[w])\otimes (\pmb{\infty},\pmb{0})-[-1]\otimes\partial_2(X_1)\\
			=&[-1]\otimes(\pmb{\infty},\pmb{0})-[-1]\otimes\partial_2(X_1)+d_2\otimes\id([w|w]\otimes(\pmb{\infty},\pmb{0})\\
			=&[-1]\otimes(-(\pmb{0},\pmb{1})+(\pmb{\infty},\pmb{1}))+d_2\otimes\id([w|w]\otimes(\pmb{\infty},\pmb{0})\\
			=&(h_1^{-1}[-1]-g_1^{-1}[-1])\otimes(\pmb{\infty},\pmb{0})+d_2\otimes\id([w|w]\otimes(\pmb{\infty},\pmb{0})\\
			=&d_2\otimes\id(\{[h^{-1}_1|-1]-[-1|h^{-1}_1]-[g^{-1}_1|-1]+[-1|g^{-1}_1]+[w|w]\}\otimes(\pmb{\infty},\pmb{0}))\\
		\end{align*}
		where $h_1=\mtxx{1}{1}{0}{1}$ and $g_1=\mtxx{0}{1}{-1}{1}$. This lifts to the element
		\begin{align*}
			\theta:&=\{[h^{-1}_1|-1]-[-1|h^{-1}_1]-[g^{-1}_1|-1]+[-1|g^{-1}_1]\}\otimes(\pmb{\infty},\pmb{0})
			\in B_2(\SL_2(A))\otimes_{\SL_2(A)} X_1(A^2).
		\end{align*}
		Applying $\id\otimes\partial_1$, we obtain
		\begin{align*}
			(\id\otimes\partial_1)(\theta)&=\{w([h^{-1}_1|-1]-[-1|h^{-1}_1]-[g^{-1}_1|-1]+[-1|g^{-1}_1])\\
			&-[h^{-1}_1|-1]+[-1|h^{-1}_1]+[g^{-1}_1|-1]-[-1|g^{-1}_1]\}\otimes(\pmb{\infty})
		\end{align*}
		in $B_2(\SL_2(A))\otimes_{\SL_2(A)} X_1(A^2)$. This represents the image $d^2_{2,1}(\chi_1)\in H_2(B(A),\z)$. For $g\in\SL_2(A)$, consider the element $c(g,-1)=([g|-1]-[-1|g])\otimes(\pmb{\infty})\in H_2(B(A),\z)$. These elements satisfy the following properties.
		\begin{enumerate}
			\item For $h\in B(A)$ and $g\in\SL_2(A)$ we have
			\[
			c(hg,-1)=c(h,-1)+c(g,-1)
			\]
			\item For any $g\in\SL_2(A)$,
			\[
			wc(g,-1)=c(wg,-1)-c(w,-1).
			\]
			\item $(w[w|w]-[w|w])\otimes(\pmb{\infty})=-c(w,-1)$.
			\item For $h\in B(A)$ we have $-c(h,-1)=c(h^{-1},-1)$.
		\end{enumerate}
		
		These identities imply
		\begin{align*}
			d^2_{2,1}(\chi_1)&=c(wh_1^{-1},-1)-c(h_1^{-1},-1)-c(wg_1^{-1},-1)+c(g_1^{-1},-1)-c(w,-1)\in H_2(B(A),\z)
		\end{align*}
		Using the identities $wg_1^{-1}=h_1^{-1}wh_1^{-1}$ and $g_1^{-1}=-h_1w$, together with the properties of $c(\cdot,\cdot)$ above, we obtain
		\[
		d^2_{2,1}(\chi_1)=-3c(h^{-1}_1,-1)=c(h_1,-1)=c\left(\mtxx{1}{1}{0}{1},-1\right).
		\]
	\end{proof}
	%\begin{lem}\label{ker-d-alpha}
	%    If $A$ is a domain with $\char(A)\neq 2$, the kernel of $d_{2,1}^2\circ\alpha:\GG_A\to H_2(B(A),\z)\simeq \aa\wedge\aa\oplus \SS_2$ have at most two elements. Moreover, if $\lan a\ran\in \ker(d_{2,1}^2\circ\alpha)$, then $\lan a \ran$ is null or $\lan a\ran=\lan \zeta\ran$ were $\zeta$ is a primitive $l$-th root of unit and $l$ is even. In particular if $-1\notin\aa^2$, $\lan a\ran=\lan -1\ran$.
	%\end{lem}
	The next lemma is a reformulation of \cite[Lemma 4.5]{B-E--2023}.
	%\begin{lem}\label{ker-d-alpha}
	%    If $A$ is a domain with $\char(A)\neq 2$. If the element $-1\wedge a=0$ in $\aa\wedge\aa$, then the quadratic square class $\lan a \ran$ is null or $\lan a\ran=\lan \zeta\ran$ were $\zeta$ is a primitive $l$-th root of unit with $l$ even. And in this case, in particular if $-1\notin\aa^2$, $\lan a\ran=\lan -1\ran$.
	%\end{lem}
	\begin{lem}\label{ker-d-alpha}
		Let $A$ be a domain with $\char(A)\neq 2$. If $-1\wedge a=0$ in $\aa\wedge\aa$, then $a=p^l\nu$, where $\nu$ is an $l$-th root of unity and $l$ is an even integer.
	\end{lem}
	\begin{proof}
		Write $\aa=\varinjlim H$, where $H$ runs over the finitely generated subgroups of $\aa$. Since the wedge product (second homology) commutes with direct limits, we may regard $-1\wedge a\in \aa\wedge \aa$ as an element of $H\wedge H$ for a suitable $H$ containing $-1$ and $a$. Decompose $H=F\oplus T$ into its free part $F$ and its torsion part $T$ (the roots of unity in $H$); note that $T$ is cyclic and of even order, since $-1\in T$. By the K\"unneth formula, we obtain
		\[
		H\wedge H=F\wedge F\oplus F\otimes T\oplus T\wedge T.
		\]
		The summand $T\wedge T$ vanishes because $T$ is finite cyclic ($A$ being a domain). Write $a=q\nu$ with $q\in F$ and $\nu\in T$; then $a\wedge (-1)=(0,q\otimes(-1),\nu\wedge (-1))=(0,q\otimes(-1),0)$. Setting $T\simeq \z/l$ with $l=|T|$, we see that if $-1\wedge a=0$, then $\displaystyle 0=q\otimes(-1)\in F\otimes T\simeq\frac{F}{lF}$, so that $q=p^l$ for some $p\in F$, with $l$ even. %Finally $\lan a\ran=\lan pw\ran=\lan w\ran$ because $l$ is even. Consider a generator $\zeta$ of $T$, then $\lan w\ran=\lan \zeta^r\ran$ and
		%\[
		%    \lan w\ran=\begin{cases}
			%        \lan 1\ran, & r \textrm{ is even}\\
			%        \lan\zeta\ran, & r \textrm{ is odd}
			%    \end{cases}.
		%\]
		%In the last case above, we have $\zeta^{l/2}=-1$ ($l=|T|$) and if $-1\notin\aa^2$ we have that $l/2$ is odd and $\lan\zeta\ran=\lan\zeta^{l/2}\ran=\lan -1\ran$.
	\end{proof}
	\begin{lem}
		\label{ker-wed-2}
		If $A$ is a domain with $\char(A)\neq 2$, then the $2$-torsion subgroup of $\bigwedge^2\aa$ is generated by the elements $q\wedge (-1)$ with $q\in\aa$.
	\end{lem}
	\begin{proof}
		Let $\theta\in\bigwedge^2\aa$. Since $\aa$ is the direct limit of its finitely generated subgroups, we may choose such a subgroup $H$ which contains $-1$ and satisfies $\theta\in\bigwedge^2H$. Let $H=L\oplus T$ be the decomposition into free and torsion parts, respectively. By the K\"unneth formula,
		\[
		H\wedge H=L\wedge L\oplus L\otimes T\oplus T\wedge T=L\wedge L\oplus L\otimes T
		\]
		where the summand $T\wedge T$ vanishes because $T$ is finite cyclic. Since $\theta$ is a torsion element of $H\wedge H$ and $L\wedge L$ is free, we have $\displaystyle\theta\in L\otimes T\simeq \frac{L}{|T|L}$. We may therefore write $\theta=p\wedge\zeta$, where $p\in L$ and $\zeta\in T$ is a generator of $T$. Since $2\theta=0$, we have $p^2=q^{|T|}$ for some $q$, and hence (note that $T$ has even order because $-1\in T$)
		\[
		p=q^{\frac{|T|}{2}}
		\]
		Finally, $p\wedge\zeta=q\wedge\zeta^{|T|/2}=q\wedge(-1)$.
	\end{proof}
	\begin{cor}\label{ker-gam}
		Let $A$ be a domain with $\char(A)\neq 2$. Then every $\theta\in\ker(\gamma)$ is of the form $\theta=K+\Lan \nu\Ran\otimes(-1)$, where $K\in\II^2_A\otimes\mu_2(A)$ and $\nu$ is an $l$-th root of unity in $A$ with $l$ even.
	\end{cor}
	\begin{proof}
		Let $x=\displaystyle\sum_{i=0}^n\Lan a_i\Ran\otimes(-1)\in\ker(\gamma)$. Applying $\gamma$ to $x$ and using Lemma \ref{alpha-gam-chi} (3), we obtain the equation
		\[
		0=\alpha(\lan\prod_i a_i\ran)+\sum_i\Lan a_i\Ran\chi_1.
		\]
		
		Applying the differential $d^2_{2,1}$, we obtain an element of $H_2(B(A),\z)$. Lemma \ref{alpha-gam-chi} (5), Lemma \ref{d2-chi} and the decomposition $H_2(B(A),\z)\simeq H_2(T(A),\z)\oplus \SS_2$ then imply that
		\[
		0=-1\wedge\prod_{i=1}^na_i=c(-1,\prod_{i=1}^nh_{a_i+1}).
		\]
		Hence, by Lemma \ref{ker-d-alpha}, $\displaystyle\lan\prod_{i=1}^na_i\ran=\lan\nu\ran$, where $\nu$ is a primitive $l$-th root of unity with $l$ even. Note that
		\[
		\Lan\nu\Ran\otimes(-1)=\Lan\prod_{i=1}^na_i\Ran\otimes(-1)=\Lan a_1\Ran\Lan\prod_{i>1}^na_i\Ran\otimes(-1)+\Lan a_1\Ran\otimes(-1)+\Lan\prod_{i>1}^na_i\Ran\otimes(-1)
		\]
		and then
		\[
		\Lan a_1\Ran\otimes(-1)+\Lan\prod_{i>1}^na_i\Ran\otimes(-1)=\Lan\nu\Ran\otimes(-1)+\Lan a_1\Ran\Lan\prod_{i>1}^na_i\Ran\otimes(-1)
		\]
		where the second summand on the right-hand side lies in $\II^2_A\otimes\mu_2(A)$. Iterating this process in order to reduce $\displaystyle\Lan\prod_{i>1}^na_i\Ran\otimes(-1)$, we obtain
		\[
		\sum_{i=1}^n\Lan a_i\Ran\otimes(-1)=K+\Lan\nu\Ran\otimes(-1)
		\]
		where $K\in \II^2_A\otimes\mu_2(A)$.
	\end{proof}
	%%%%%%%%%%%%%%%%%%%%%%%%%%%%%%%%%%%%%%%%%%%%
	\section{The element \texorpdfstring{$\theta_a$}{Oa}}
	%%%%%%%%%%%%%%%%%%%%%%%%%%%%%%%%%%%%%%%%%%%%
	For $a\in\aa$, consider the element $\theta_a=[w]\otimes\partial_2(X_{-a^{-1}}-X_a)\in E^2_{2,1}$. Then
	\begin{equation}
		\label{tht-id}
		\begin{split}
			\theta_a&=[w]\otimes\partial_2(X_{-a^{-1}}+X'_{-a^{-1}}-X'_{-a^{-1}}-X_a)\\
			&=[w]\otimes Y - [w]\otimes\partial_2(X'_{-a^{-1}}+X_a)\\
			&=[w]\otimes Y - [w]\otimes\partial_2(wX_a+X_a)\\
			&=[w]\otimes Y - \{w[w]+[w]\}\otimes\partial_2(X_a)\\
			&=[w]\otimes Y - [-1]\otimes\partial_2(X_a)
		\end{split}
	\end{equation}
	\begin{prp}
		\label{sqr-plus-tht}
		Let $a\in\aa$ be such that $a^2+1\in\aa$. Then $\Lan a^2+1\Ran\theta_a=0$.
	\end{prp}
	\begin{proof}
		For $a\in\aa$, consider the elements $\tilde{X}_a=(\pmb{a},\pmb{\infty},\pmb{0})$ and $\tilde{X}'_a=(\pmb{a},\pmb{0},\pmb{\infty})$ of $X_2(A^2)$. The element $\partial_3((\pmb{\infty},\pmb{0},\pmb{-a^{-1}},\pmb{a}))\in Z_2(A^2)$ makes sense precisely because $a^2+1\in\aa$, and hence $\theta_a=[w]\otimes\partial_2((\pmb{0},\pmb{-a^{-1}},\pmb{a})-(\pmb{\infty},\pmb{-a^{-1}},\pmb{a}))$, and setting $\rho_a=\mtxx{a}{\frac{1}{a^2+1}}{-1}{\frac{a}{a^2+1}}\in\SL_2(A)$, we have $\theta_a=[w]\otimes\rho_a\partial_2(\tilde{X}_{-a(a^2+1)}-\tilde{X}_{\frac{a^2+1}{a}})$. It is not difficult to show that $\rho_a^{-1}w\rho_a=w(a^2+1)$, and therefore
		\begin{align*}
			\theta_a&=[\rho_a^{-1}w\rho_a]\otimes\rho_a^{-1}\rho_a\partial_2(\tilde{X}_{-a(a^2+1)}-\tilde{X}_{\frac{a^2+1}{a}})\\
			&=[w(a^2+1)]\otimes\partial_2(\tilde{X}_{-a(a^2+1)}-\tilde{X}_{\frac{a^2+1}{a}})\\
			&=\lan -(a^2+1)\ran(-[-w]\otimes\partial_2(\tilde{X}_{-a^{-1}}-\tilde{X}_{a}))\\
			&=\lan -(a^2+1)\ran ([w]\otimes\partial_2(\tilde{X}_{-a^{-1}}-\tilde{X}_a))\\
			&=\lan-(a^2+1)\ran\tilde{\theta}_a
		\end{align*}
		where $\tilde{\theta}_a:=[w]\otimes\partial_2(\tilde{X}_{-a^{-1}}-\tilde{X}_a)$. Now, using \eqref{tht-id} and the matrix $\mtxx{1}{-a^{-1}}{a}{0}\in\SL_2(A)$, which sends $X_a$ to $\tilde{X}_a$, we obtain
		\begin{align*}
			\theta_a&=[w]\otimes Y-[-1]\otimes\partial_2(X_a)\\
			&=[w]\otimes Y-[-1]\otimes\partial_2(\tilde{X}_a)\\
			&=[w]\otimes\partial_2(\tilde{X}_{-a^{-1}}+\tilde{X}'_{-a^{-1}})-[-1]\otimes\partial_2(\tilde{X}_a)\\
			&=[w]\otimes\partial_2(\tilde{X}_{-a^{-1}}+\tilde{X}'_{-a^{-1}})-\{w[w]+[w]\}\otimes\partial_2(\tilde{X}_a)\\
			&=[w]\otimes\partial_2(\tilde{X}_{-a^{-1}}+\tilde{X}'_{-a^{-1}})-[w]\otimes\partial_2(\tilde{X}'_{-a^{-1}}+\tilde{X}_a)\\
			&=[w]\otimes\partial_2(\tilde{X}_{-a^{-1}}-\tilde{X}_a)=\tilde{\theta}_a\\
		\end{align*}
		so that $\Lan-(a^2+1)\Ran\theta_a=0$.
		Now, by Lemma \ref{alpha-gam-chi} (2) and by \eqref{tht-id}, we have
		\[
		\chi_1-\theta_a=\gamma(\Lan a\Ran\otimes(-1))
		\]
		Multiplying by $\Lan -1\Ran$ and using Lemma \ref{alpha-gam-chi} (6), we get
		\begin{align*}
			-\Lan -1\Ran\theta_a&=\Lan a\Ran\gamma(\Lan -1\Ran\otimes(-1))\\
			&=\Lan a\Ran\alpha(\lan -1\ran)\\
			&=0
		\end{align*}
		and finally
		\[
		0=\Lan -(a^2+1)\Ran\theta_a=\Lan a^2+1\Ran\Lan -1\Ran\theta_a+\Lan -1\Ran\theta_a+\Lan a^2+1\Ran\theta_a=\Lan a^2+1\Ran\theta_a.
		\]
	\end{proof}
	\begin{cor}
		\label{sqr-plus-chi}
		Let $a\in\aa$ be such that $a^2+1\in\aa$. Then $\Lan a^2+1\Ran\chi_1=0$.
	\end{cor}
	\begin{proof}
		By Lemma \ref{alpha-gam-chi} (2), we have
		\[
		\chi_a=\lan a\ran\chi_1=[wa]\otimes Y-[-1]\otimes\partial_2(X_a)=[a]\otimes Y+[w]\otimes Y-[-1]\otimes\partial_2(X_a)]
		\]
		and using \eqref{tht-id} we have
		\[
		\chi_a-\theta_a=\alpha(\lan a\ran)
		\]
		Multiplying by $\lan a\ran\Lan a^2+1\Ran$ and using Proposition \ref{sqr-plus-tht}, we obtain the required identity.
	\end{proof}
	%%%%%%%%%%%%%%%%%%%%%%%%%%%%%%%%%%%%%%%%%%%%%%%%%%%%%%%%%
	\section{The map \texorpdfstring{$H_\bullet(B(A),\z)\to H_\bullet(T(A),\z)$}{H.(B)->H.(T)}}\label{sec-hb-ht}
	%%%%%%%%%%%%%%%%%%%%%%%%%%%%%%%%%%%%%%%%%%%%%%%%%%%%%%%%%
	Consider the group extension
	\[
	\begin{tikzcd}
		0\ar[r] & N(A)\ar[r] & B(A) \ar[r] & T(A)\ar[r] & 0
	\end{tikzcd}
	\]
	where the map $B(A)\to T(A)$ sends $\mtxx{a}{b}{0}{a^{-1}}$ to $\mtxx{a}{0}{0}{a^{-1}}$, and the kernel $N(A)$ consists of the matrices $\mtxx{1}{b}{0}{1}$ with $b\in A$. Clearly $N(A)\simeq A$ and $T(A)\simeq\aa$, and the resulting action of $\aa$ on $A$ is given by $a\cdot b=a^2b$ (it is induced by conjugation).
	The map above is split by the inclusion $T(A)\to B(A)$, which yields a canonical decomposition, for any $n\geq 0$,
	\[
	H_n(B(A),\z)\simeq H_n(T(A),\z)\oplus H_n(B(A),T(A);\z),
	\]
	where $H_n(B(A),T(A);\z)$ denotes the relative homology of the pair $(B(A),T(A))$ (for definitions and properties, see \cite[page 153]{knudson2001}). This relative homology group will often be denoted by $\SS_n$.
	\begin{lem}
		\label{ht-hb-l1}
		If $A$ is a subring of $\q$, then for any $n\geq 0$
		\[ 
		H_n(B(A),\z)\simeq H_n(T(A),\z)\oplus H_{n-1}(\aa,A)
		\]
		and if $6\in\aa$ then $H_n(B(A),\z)\simeq H_n(T(A),\z)$ for $n\geq 0$.
	\end{lem}
	\begin{proof}
		See \cite[Lemma 3.5]{B-E--2023}.
	\end{proof}
	In the case of local rings (or fields) we have
	\begin{prp}\label{ht-hb-local}
		Let $A$ be a local ring with maximal ideal $\mmm_A$.
		\begin{enumerate}
			\item If $A/\mmm_A$ is infinite, or if $|A/\mmm_A|=p^d$ with $(p-1)d>2(n+1)$, then
			\[
			H_n(B(A),\z)\simeq H_n(T(A),\z)
			\]
			\item If $A$ is a domain, and if $A/\mmm_A$ is infinite or $|A/\mmm_A|=p^d$ with $(p-1)d>2n$, then
			\[
			H_n(B(A),\z)\simeq H_n(T(A),\z)
			\] 
		\end{enumerate}
	\end{prp}
	\begin{proof}
		See \cite[Proposition 3.19 ]{hutchinson2017}.
	\end{proof}
	\begin{exm}\label{exa-hb-ht}
		\begin{enumerate}
			\item For the ring $\z$, by Lemma \ref{ht-hb-l1}, $\SS_3=H_2(\{\pm 1\},\z)=0$. Then
			\[
			H_3(B(\z),\z)\simeq H_3(T(\z),\z)
			\]
			\item For a finite field $\F_q$ with $q=p^d$, if $(p-1)d>6$, then Proposition \ref{ht-hb-local} gives
			\[
			H_3(B(\F_q),\z)\simeq H_3(T(\F_q),\z)
			\]
			\item For the finite local rings $\z/p^r$, the residue field is $\z/p$; hence, if $p>7$, Proposition \ref{ht-hb-local} gives
			\[
			H_3(B(\z/p^r),\z)\simeq H_3(T(\z/p^r),\z)
			\]
		\end{enumerate}
	\end{exm}
	For $m\in\N$, we denote by $A_m$ the subring $\z\left[\frac{1}{m}\right]$ of $\q$.
	\begin{lem}
		Let $p$ be a prime number. Then
		\begin{enumerate}
			\item $H_1(B(A_p),\z)\simeq H_1(T(A_p),\z)\oplus\z/(p^2-1)$.
			\item For any $n\geq 2$, $H_n(B(A_2),\z)\simeq H_n(T(A_2),\z)$.
			\item For $p\neq 2$ and $n\geq 2$, we have $H_n(B(A_p),\z)\simeq H_n(T(A_p),\z)\oplus\z/2$.
		\end{enumerate}
	\end{lem}
	\begin{proof}
		See \cite[Lemma 3.7]{B-E--2023}.
	\end{proof}
	A similar technique extends the above result to the case where $m$ has two or more prime factors; we treat the cases in which $2\in\aa$ but $3\notin\aa$.
	\begin{prp}\label{prp-hb-ht-am}
		Let $m\in\N$, $r\geq 1$ and $m=2^\alpha\cdot p_1^{\alpha_1}\cdots p_r^{\alpha_r}$ with $p_i\neq 3$ for $i\leq r$.
		\begin{enumerate}
			\item $H_1(B(A_m),\z)\simeq H_1(T(A_m),\z)\oplus \z/3$.
			\item $H_2(B(A_m),\z)\simeq H_2(T(A_m),\z)\oplus(\z/3)^r$.
			\item $H_3(B(A_m),\z)\simeq H_3(T(A_m),\z)\oplus(\z/3)^{\frac{(r-1)r}{2}}$.
			\item If $r=1$ then $H_n(B(A_m),\z)\simeq H_n(T(A_m),\z)$ for any $n\geq 3$.
		\end{enumerate}
	\end{prp}
	\begin{proof}
		By Lemma \ref{ht-hb-l1}, it suffices to compute $H_{k}(A_m^\times,A_m)$ for $k=0,1,2$. Since $A_m^\times=\lan-1,p_1,\dots,p_r\ran$, consider the split extension
		\[
		\begin{tikzcd}
			0\to \lan 2,p_1,\dots,p_r\ran\ar[r] & A_m^\times \ar[r] & \mu_2(A_m)\ar[r] & 0 
		\end{tikzcd}
		\]
		this yields a Lyndon/Hochschild-Serre spectral sequence
		\[
		\prescript{I}{}{E}^2_{p,q}=H_p(\lan 2,p_1,\dots,p_r\ran,H_q(\mu_2(A_m),A_m))\implies H_{p+q}(A_m^\times,A_m).
		\]
		
		Since $\mu_2(A_m)$ acts trivially on $A_m$, we have $H_0(\mu_2(A_m),A_m)=A_m$, and the $2$-divisibility of $A_m$ implies that $H_q(\mu_2(A_m),A_m)=0$ for $q\geq 1$. Hence, by the convergence of the spectral sequence, $H_p(A_m^\times,A_m)\simeq\prescript{I}{}{E}^2_{p,0}=H_p(\lan 2,p_1,\dots,p_r\ran,A_m)$.
		
		Consider now the split extension
		\[
		\begin{tikzcd}
			0\to \lan p_1,\dots,p_r\ran\ar[r] & \lan2,p_1,\dots,p_r\ran \ar[r] & \lan 2\ran\ar[r] & 0 
		\end{tikzcd}
		\]
		this yields a spectral sequence
		\[
		\prescript{II}{}{E}^2_{p,q}=H_p(\lan p_1,\dots,p_r\ran,H_q(\lan 2\ran,A_m))\implies H_{p+q}(\lan2,p_1,\dots,p_r\ran,A_m).
		\]
		
		Now, $H_0(\lan 2\ran,A_m)\simeq A_m/(2^2-1)\simeq\z/3$ because $3\nmid m$. Since $\lan 2\ran$ is an infinite cyclic group and $A_m$ is torsion free, we have $H_q(\lan 2\ran,A_m)=0$ for $q\geq 1$. Hence, by the convergence of this second spectral sequence, $H_p(A_m^\times,A_m)\simeq H_p(\lan 2,p_1,\dots,p_r\ran,A_m)\simeq\prescript{II}{}{E}^\infty_{p,0}=\prescript{II}{}{E}^2_{p,0}=H_p(\lan p_1,\dots,p_r\ran,\z/3)$.
		
		Finally, $\prescript{II}{}{E}^2_{p,0}=H_p(\lan p_1,\dots,p_r\ran,\z/3)\simeq H_p(\lan p_1,\dots,p_r\ran,\z)\otimes \z/3$, because $\lan p_1,\dots,p_r\ran$ acts trivially on $\z/3$ (here we also use the universal coefficient theorem \cite[III.1, Exercise 3]{brown1994}). For $p=0$ we obtain $\prescript{II}{}{E}^2_{0,0}=\z/3$, which proves (1). For $p\geq 1$ we distinguish two cases.
		\begin{enumerate}
			\item Suppose $r>1$. For $p=1$ we have $\prescript{II}{}{E}^2_{1,0}=\z^r\otimes(\z/3)\simeq\left(\z/3\right)^r$, and for $p=2$ we have $\prescript{II}{}{E}^2_{2,0}=\z^{\frac{(r-1)r}{2}}\otimes(\z/3)\simeq\left(\z/3\right)^{\frac{(r-1)r}{2}}$ (by the K\"unneth formula); this proves (2) and (3).
			\item Suppose $r=1$. The group $\lan p_1\ran$ is infinite cyclic, so $\prescript{II}{}{E}^2_{p,0}=0$ for $p\geq 2$, while for $p=1$ we have $\prescript{II}{}{E}^2_{1,0}=\z\otimes\z/3\simeq\z/3$; this proves (4).
		\end{enumerate}
	\end{proof}
	%%%%%%%%%%%%%%%%%%%%%%%%%%%%%%%%%%%%%%%%%%%%%%%%%%%%%%%%%
	\section{The homology \texorpdfstring{$H_3(\SL_2(A),\z)$ and $H_3(\SL_2(A),\SM_2(A);\z)$}{H3(SL2(A),z) and H3(SL2(A),SM2(A);z)}}
	%%%%%%%%%%%%%%%%%%%%%%%%%%%%%%%%%%%%%%%%%%%%%%%%%%%%%%%%%
	Let $A$ be a ring and consider the following commutative diagram with exact rows.
	\begin{equation}
		\begin{tikzcd}
			& \hat{E}^2_{2,1}\ar[r,"\alpha"]\ar[d,twoheadrightarrow,"\hat{d}^2_{2,1}"] & E^2_{2,1}\ar[r]\ar[d,"d^2_{2,1}"] & \coker(\alpha) \ar[r] \ar[d,"d'"] & 0\\
			0 \ar[r] & \{\pm 1\}\wedge\aa\ar[r] & \aa\wedge\aa\oplus\SS_2 \ar[r,twoheadrightarrow] & \frac{\aa\wedge\aa}{\aa\wedge\{\pm 1\}}\oplus\SS_2 \ar[r]& 0 
		\end{tikzcd}
	\end{equation}
	
	Here $d'$ is induced by $d^2_{2,1}$. Setting $\widetilde{\coker(\alpha)}:=\ker(d')$, the snake lemma yields the exact sequence
	\begin{equation}
		\label{alph-inf}
		\begin{tikzcd}
			\hat{E}^{\infty}_{2,1}\ar[r,"\alpha"] & E^{\infty}_{2,1}\ar[r] & \widetilde{\coker(\alpha)} \ar[r] & 0
		\end{tikzcd}
	\end{equation}
	
	The next proposition relates the cokernel of $\alpha$ to the cokernel of the map induced by the inclusion $H_3(\SM_2(A),\z)\to H_3(\SL_2(A),\z)$. 
	\begin{prp}\label{prp-sm2-sl2}
		Let $A$ be a ring satisfying the following conditions:
		\begin{enumerate}
			\item The complex $X_{\bullet}(A^2)\to\z$ is exact in dimension $<2$.
			\item $H_3(B(A),\z)\simeq H_3(T(A),\z)$.
		\end{enumerate}
		Then there is an exact sequence
		\[
		H_3(\SM_2(A),\z)\to H_3(\SL_2(A),\z)\to \widetilde{\coker(\alpha)}\to 0.
		\]
	\end{prp}
	\begin{proof}
		The proof is the same as that of Theorem 6.6 in \cite{B-E--2023}. Analyzing the map \eqref{map-s-seq} of spectral sequences $\hat{E}^1_{p,q}\overset{\inc}{\longrightarrow} E^1_{p,q}$, we obtain filtrations (the equalities on the right follow from $\hat{E}^{\infty}_{3,0}=E^{\infty}_{3,0}=0$)
		\[
		\begin{tikzcd}
			0\ar[r,hook] & \hat{F}_0 \ar[r,hook]\ar[d] & \hat{F}_1 \ar[r,hook]\ar[d] & \hat{F}_2 \ar[r,equal]\ar[d] & \hat{F}_3 = H_3(\SM_2(A),\z)\ar[d]\\
			0\ar[r,hook] & F_0 \ar[r,hook] & F_1 \ar[r,hook] & F_2 \ar[r,equal] & F_3 = H_3(\SL_2(A),\z)
		\end{tikzcd}
		\]
		This yields commutative diagrams with exact rows
		\begin{equation}
			\label{rel_2-1}
			\begin{tikzcd}
				0 \ar[r] & \hat{F}_1 \ar[r]\ar[d] & H_3(\SM_2(A),\z)\ar[r]\ar[d] & \hat{E}^{\infty}_{2,1}\ar[r]\ar[d] & 0\\
				0 \ar[r] & F_1 \ar[r] & H_3(\SL_2(A),\z)\ar[r] & E^{\infty}_{2,1}\ar[r] & 0.\\
			\end{tikzcd}    
		\end{equation}
		and
		\begin{equation}
			\label{rel_1-2}
			\begin{tikzcd}
				0 \ar[r] & \hat{E}^{\infty}_{0,3}=\hat{F}_0 \ar[r]\ar[d] & \hat{F}_1 \ar[r]\ar[d] & \hat{E}^{\infty}_{1,2}\ar[r]\ar[d] & 0\\
				0 \ar[r] & E^{\infty}_{0,3}=F_0 \ar[r] & F_1 \ar[r] & E^{\infty}_{1,2}\ar[r] & 0.\\
			\end{tikzcd}    
		\end{equation}
		Applying the snake lemma to diagram \eqref{rel_2-1}, we see that it suffices to prove that the map $\hat{F}_1\to F_1$ is surjective. To this end, observe that in diagram \eqref{rel_1-2} the groups $\hat{E}^{\infty}_{1,2}$ and $E^{\infty}_{1,2}$ are quotients of $H_2(T(A),\z)$, which implies that the right-hand vertical map is surjective. Moreover, the condition $H_3(B(A),\z)\simeq H_3(T(A),\z)$ implies, by the same argument, that the left-hand vertical map is surjective. Therefore $\hat{F}_1\to F_1$ is surjective.
	\end{proof}
	Our first main theorem analyzes $\coker(\alpha)$ and provides weaker conditions under which an exact sequence as in \cite[Theorem 6.6]{B-E--2023} holds. Consider the following completion of diagram \eqref{dia-gam-alpha}.
	\begin{equation}\label{dia-gam-alpha-2}
		\begin{tikzcd}
			& & \II_A\otimes\mu_2(A)\ar[d,"\gamma"] & & \\
			& \hat{E}^2_{2,1}\simeq\GG_A \ar[r,"\alpha:=\inc_*"]\ar[d,twoheadrightarrow,"\delta'':=\delta\circ\alpha"] & E^2_{2,1} \ar[r,"\pi"]\ar[d,twoheadrightarrow,"\delta"] & \coker(\alpha)\ar[r]\ar[d,twoheadrightarrow,"\delta'"] & 0\\
			0 \ar[r] & \II_A\psi_1(-1)\ar[r] & \RP_1(A)\ar[r] & \frac{\RP_1(A)}{\II_A\psi_1(-1)}\ar[r] & 0
		\end{tikzcd}
	\end{equation}
	The surjectivity of the left-hand vertical map follows from Lemma \ref{alpha-gam-chi} (4), and the surjectivity of the right-hand vertical map follows from the five lemma.
	\begin{thm}\label{thm-main}
		Let $A$ be a ring such that $X_{\bullet}(A^2)\to\z$ is exact in dimension $<2$, and consider the maps of diagram \eqref{dia-gam-alpha-2}. Then:
		\begin{enumerate}
			\item There exists an exact sequence
			\[
			\II_A\otimes\mu_2(A)\to\coker(\alpha)\to\frac{\RP_1(A)}{\II_A\psi_1(-1)}\to 0
			\]
			\item $\im(\gamma)\subseteq\im(\alpha)$ if and only if $\delta'$ is an isomorphism.
			\item $\im(\alpha)\subseteq\im(\gamma)$ if and only if $\II_A\psi_1(-1)=0$.
			\item If $A$ is a domain, then $\im(\gamma)\subseteq\im(\alpha)$ if and only if $\im(\alpha)\subseteq\im(\gamma)$ and $\II_A^2\otimes\mu_2(A)\subseteq\ker(\gamma)$.
		\end{enumerate}
	\end{thm}
	\begin{proof}
		Applying the snake lemma to diagram \eqref{dia-gam-alpha-2}, we obtain the exact sequence
		\begin{equation}\label{snk-dia-gam-alpha-2}
			\begin{tikzcd}
				\ker(\delta'')\ar[r,"\bar{\alpha}"] & \ker(\delta)=\im(\gamma)\ar[r,"\bar{\pi}"] & \ker(\delta') \ar[r] & 0.
			\end{tikzcd}
		\end{equation}
		For (1), the left-hand map is the composite $\pi\circ\gamma$ and the right-hand map is $\delta'$; exactness follows from the exact sequence above.
		
		For (2), suppose that $\im(\gamma)\subseteq\im(\alpha)$, and let $x=\gamma(y)=\alpha(z)$ with $y\in\II_A\otimes\mu_2(A)$ and $z\in\hat{E}^2_{2,1}$. Then $\delta''(z)=\delta(\alpha(z))=0$, so $\bar{\alpha}$ is surjective and $\ker(\delta')=0$. The converse follows from the exact sequence \eqref{snk-dia-gam-alpha-2}.
		\\
		For (3), it is clear that if $\im(\alpha)\subseteq\im(\gamma)$, then $\delta''=\delta\circ\alpha=0$, which implies that $\II_A\psi_1(-1)=0$. The converse follows from diagram \eqref{dia-gam-alpha-2}.
		\\
		For (4), let $A$ be a domain and assume that $\im(\gamma)\subseteq\im(\alpha)$; it is then clear that $\II_A^2\otimes\mu_2(A)\subseteq \ker(\gamma)$. We shall only treat the case $-1\notin\aa^2$, since the remaining case follows from \cite[Corollary 6.4]{B-E--2023}; similarly, we may assume that $\char(A)\neq 2$, the other case being obvious. Let $a\in\aa$. By the hypothesis and Lemma \ref{alpha-gam-chi} (3), there exists $b\in\aa$ such that 
		\[
		\gamma(\Lan a\Ran\otimes(-1))=\alpha(\lan b\ran)=\alpha(\lan a\ran)+\Lan a\Ran\chi_1
		\]
		Applying $d^2_{2,1}$ to the right-hand equality and using Lemma \ref{ker-d-alpha}, we obtain $\lan b\ran=\lan a\ran$ or $\lan b\ran=\lan -a\ran$ (recall the decomposition $H_2(B(A),\z)=\aa\wedge\aa\oplus \SS_2$). In the first case we are done. In the second case, the equation above and Lemma \ref{alpha-gam-chi} (6) imply that $\Lan a\Ran\chi_1=\alpha(\lan -1\ran)=\gamma(\Lan -1\Ran\otimes(-1))$, and finally $\alpha(\lan a\ran)=\gamma(\Lan a\Ran\otimes(-1)-\Lan -1\Ran\otimes (-1))$. 
		
		Conversely, take an element $\sum\Lan a_i\Ran\otimes(-1)$. Lemma \ref{alpha-gam-chi} (3) and the hypotheses imply that, for some $b\in\aa$,
		\begin{equation}
			\label{eq-prf4}
			\alpha(\lan\prod_{i}a_i\ran)=\gamma(\Lan b\Ran\otimes(-1)).
		\end{equation}
		
		By Lemma \ref{alpha-gam-chi} (3), we have $\gamma(\Lan b\Ran\otimes(-1))=\alpha(\lan b\ran)+\Lan b\Ran\chi_1$. Applying $d^2_{2,1}$ to the equation $\alpha(\lan\prod_{i}a_i\ran)=\alpha(\lan b\ran)+\Lan b\Ran\chi_1$ and using the decomposition $H_2(B(A),\z)=\aa\wedge\aa\oplus \SS_2$, we obtain
		\[
		(b\cdot\prod_ia_i)\wedge (-1)=0 \text{ and } \Lan b\Ran c(h_1,-1)=0,
		\]
		and Lemma \ref{ker-d-alpha} (note that $-1\notin\aa^2$) leads to two cases:
		\[
		\lan b\ran=\lan\prod a_i\ran\text{ or } \lan b\ran=\lan-\prod a_i\ran.
		\]
		
		Now, the hypothesis $\II_A^2\otimes\mu_2(A)\subseteq\ker(\gamma)$ implies, for any $x,y\in\aa$, that
		\[
		\gamma(\Lan xy\Ran\otimes(-1))=\gamma(\Lan x\Ran\otimes(-1)+\Lan y\Ran\otimes(-1)).
		\]
		Hence, in the first case we are done by substituting $b$ in equation \eqref{eq-prf4} and using the identity above. In the second case, the same substitution yields $\alpha(\lan\prod_{i}a_i\ran)=\gamma(\Lan -\prod a_i\Ran\otimes(-1))=\gamma(\sum\Lan a_i\Ran\otimes(-1))+\gamma(\Lan -1\Ran\otimes(-1))$. Using Lemma \ref{alpha-gam-chi} (6), we obtain $\alpha(\lan-\prod_{i}a_i\ran)=\gamma(\sum\Lan a_i\Ran\otimes(-1))$. 
	\end{proof}
	\begin{exm}\label{ex-im-eq}
		\begin{enumerate}
			\item The Euclidean domain $\z$ is a universal $\GE_2$-ring which satisfies $\im(\gamma)\subseteq\im(\alpha)$ by Lemma \ref{alpha-gam-chi} (3) and (7). Hence, by Theorem \ref{thm-main} (4), $\im(\gamma)=\im(\alpha)$.
			\item It is known that a finite field $\F_q$ with $q=p^r$ and $p>2$ has exactly two square classes, $\{1,\langle u\rangle\}$, where $u$ is a non-square. When $-1$ is not a square (the other case was treated in \cite{B-E--2023}), we have $\im(\gamma)=\im(\alpha)$ by Lemma \ref{alpha-gam-chi} (7).
			\item The rings $\z/p^r$ are local, hence universal $\GE_2$-rings; moreover, $(\z/p^r)^\times\simeq\z/p^{r-1}(p-1)$ and the set of square classes is $\{1,\langle u\rangle\}$, where $u$ is a non-square. If $-1$ is not a square, then $\im(\gamma)\subseteq\im(\alpha)$ by Lemma \ref{alpha-gam-chi}. At the same time, $\II_A\psi_1(-1)=0$, and therefore $\im(\gamma)=\im(\alpha)$.
			\item Consider the domains $\z\left[\frac{1}{2}\right]$, $\z\left[\frac{1}{10}\right]$ and $\z\left[\frac{1}{290}\right]$, which are universal $\GE_2$-rings, and observe that $5=2^2+1$ and $29=5^2+2^2=2^2\{(5/2)^2+1\}$. Using the identity $\Lan ab\Ran\chi_1=\Lan a\Ran\Lan b\Ran\chi_1+\Lan a\Ran\chi_1+\Lan b\Ran\chi_1$, together with Corollary \ref{sqr-plus-chi} and Lemma \ref{alpha-gam-chi} (7), we see that every $a\in\aa$ in these rings satisfies $\Lan a\Ran\chi_1=0$; by Lemma \ref{alpha-gam-chi} (3), this gives the inclusion $\im(\gamma)\subseteq\im(\alpha)$, and Theorem \ref{thm-main} (4) implies $\im(\gamma)=\im(\alpha)$.
			\item The ring of Eisenstein integers $\z[\zeta]$, where $\zeta^2+\zeta+1=0$, is a universal $\GE_2$-ring. Its unit group $\z[\zeta]^\times=\{\pm 1,\pm\zeta,\pm\zeta^2\}$ is cyclic of order $6$, so that $\GG_{\z[\zeta]}=\{1,\lan-1\ran\}$. Lemma \ref{alpha-gam-chi} (3) and (7), together with Theorem \ref{thm-main} (4), imply that $\im(\gamma)=\im(\alpha)$.
		\end{enumerate}
	\end{exm}
	Our second main theorem establishes the exact sequence \eqref{sm2-sl2-seq} when $\im(\gamma)=\im(\alpha)$. First, note that if $\im(\gamma)\subseteq\im(\alpha)$, then Lemma \ref{alpha-gam-chi} (5) yields the commutativity of the following diagram.
	\begin{equation}
		\begin{tikzcd}
			& \II_A\otimes\mu_2(A)\ar[r,"\gamma"]\ar[d,"\Delta"] & E^2_{2,1}\ar[r,"\delta"]\ar[d,"d^2_{2,1}"] & \RP_1(A)\ar[r]\ar[d,"\Delta'"] & 0\\
			0\ar[r] & \aa\wedge\mu_2(A)\ar[r] & \aa\wedge\aa\oplus \SS_2\ar[r] & \frac{\aa\wedge\aa}{\aa\wedge\mu_2(A)}\oplus \SS_2\ar[r] & 0
		\end{tikzcd}.
	\end{equation}
	\begin{lem}\label{ker-rba}
		If $\im(\gamma)\subseteq\im(\alpha)$, then $\RB(A)=\ker(\Delta')$.
	\end{lem}
	\begin{proof}
		First, $d^2_{2,1}\circ\gamma=d^2_{2,1}(X^\tau)$, and if $\im(\gamma)\subseteq\im(\alpha)$, then $E^3_{0,3}(\SL_2(A),X^\tau)=\frac{\aa\wedge\aa}{\aa\wedge\mu_2(A)}\oplus\SS_2$ by Lemma \ref{alpha-gam-chi} (5). It is then straightforward to check that $\Delta'=d^3_{2,1}(X^\tau)$ (the proof is the same as that of Lemma 2.3 in \cite{B-E-2025}).
	\end{proof}
	%\begin{thm}\label{thm-main-1}
	%    If $A$ is a ring that satisfies 
	%    \begin{enumerate}
		%        \item The complex $X_{\bullet}(A^2)\to\z$ is exact in dimension $<2$.
		%        \item $H_3(B(A),\z)\simeq H_3(T(A),\z)$.
		%    \end{enumerate}
	%    additionally, if $\im(\gamma)\subseteq\im(\alpha)$ then we have the exact sequence
	%    \begin{equation}
		%        H_3(\SM_2(A),\z)\to H_3(\SL_2(A),\z)\to \frac{\RB(A)}{\II_A\psi_1(-1)}\to 0
		%    \end{equation}
	%    and if $\im(\gamma)=\im(\alpha)$, then we have
	%    \begin{equation}
		%        H_3(\SM_2(A),\z)\to H_3(\SL_2(A),\z)\to \RB(A)\to 0
		%    \end{equation}
	%\end{thm}
	\begin{thm}
		\label{thm-main-1}
		Let $A$ be a ring satisfying the following conditions:
		\begin{enumerate}
			\item the complex $X_{\bullet}(A^2)\to\z$ is exact in dimension $<2$;
			\item $H_3(B(A),\z)\simeq H_3(T(A),\z)$.
		\end{enumerate}
		
		If, in addition, $\im(\gamma)\subseteq\im(\alpha)$, then we have the exact sequence
		\[
		H_3(\SM_2(A),\z)\to H_3(\SL_2(A),\z)\to \frac{\RB(A)}{\II_A\psi_1(-1)}\to 0
		\]
		and if $\im(\gamma)=\im(\alpha)$, then we have
		\[
		H_3(\SM_2(A),\z)\to H_3(\SL_2(A),\z)\to \RB(A)\to 0.
		\]
	\end{thm}
	\begin{proof}
		Consider the following completion of diagram \eqref{dia-gam-alpha-2}.
		\[
		{\small
			\begin{tikzcd}
				& & & \aa\wedge\mu_2(A)\ar[dd,hook] & &\\
				& & \II_A\otimes\mu_2(A)\ar[ru,"\Delta''"]\ar[dd,"\gamma" near end] & & &\\
				& \aa\wedge\mu_2(A)\ar[rr,hook] & & \aa\wedge\aa\oplus\SS_2\ar[rr,twoheadrightarrow,"\pi''"]\ar[dd,twoheadrightarrow] & & \frac{\aa\wedge\aa}{\aa\wedge\mu_(A)}\oplus\SS_2\ar[dd,equal]\\
				\hat{E}^2_{2,1}\ar[ru,twoheadrightarrow,"\hat{d}^2_{2,1}"]\ar[rr,"\alpha"]\ar[dd,twoheadrightarrow,"\delta''"] & & E^2_{2,1}\ar[ru,"d^2_{2,1}"]\ar[dd,twoheadrightarrow,"\delta"],\ar[rr,"\pi'"]\ar[ru]\ar[dd,twoheadrightarrow] & & \coker(\alpha)\ar[ru,"d'"]\ar[dd,twoheadrightarrow,"\delta'" near end] &\\
				& & & \frac{\aa\wedge\aa}{\aa\wedge\mu_(A)}\oplus\SS_2\ar[rr,equal] & & \frac{\aa\wedge\aa}{\aa\wedge\mu_(A)}\oplus\SS_2\\
				\II_A\psi_1(-1)\ar[rr,hook] & & \RP_1(A)\ar[rr,twoheadrightarrow,"\pi"]\ar[ru,"\Delta'"] & & \frac{\RP_1(A)}{\II_A\psi_1(-1)}\ar[ru] &\\
			\end{tikzcd}
		}
		\]
		where $\pi$, $\pi'$ and $\pi''$ are the canonical maps. Since $\im(\gamma)\subseteq\im(\alpha)$, the map $\delta'$ is an isomorphism, and it is not difficult to show (using the diagram above) that the map
		\[
		\bar{\Delta'}:\frac{\RP_1(A)}{\II_A\psi_1(-1)}\to\frac{\aa\wedge\aa}{\aa\wedge\mu_2(A)}\oplus\SS_2
		\]
		induced by $\Delta'$ coincides with $d'\circ(\delta')^{-1}:\frac{\RP_1(A)}{\II_A\psi_1(-1)}\to\frac{\aa\wedge\aa}{\aa\wedge\mu_2(A)}\oplus\SS_2$. Its kernel is exactly $\frac{\RB(A)}{\II_A\psi_1(-1)}$ by Lemma \ref{ker-rba}, and this group is isomorphic to $\widetilde{\coker(\alpha)}$ via $\delta'$.
		The second assertion follows from Theorem \ref{thm-main} (3).
	\end{proof}
	\begin{exm}
		\begin{enumerate}
			\item Coronado and Hutchinson studied the cases $A=\z$ and $A=\z\half$ in \cite[\S 8]{C-H2022}; our results merely confirm their findings (see Remark \ref{rem-z12}). 
			\item For the ring $A=\z\mth{10}$ we have the exact sequence
			\[
			H_3(\SM_2(\z\mth{10}),\z)\to H_3(\SL_2(\z\mth{10}),\z)\to \RB(\z\mth{10})\to 0.
			\]
			\item A finite field $\F_q$ with $q=p^d$ and $p>2$ satisfies $-1\notin \aa^2$ when $q\equiv 3 \pmod 4$. In this case, if $(p-1)d>6$, we have 
			\[
			H_3(\SM_2(\F_q),\z)\to H_3(\SL_2(\F_q),\z)\to \RB(\F_q)\to 0
			\]
			\item For the finite local rings $\z/p^r$ with $p>9$ and $-1\notin\aa^2$, we have
			\[
			H_3(\SM_2(\z/p^r),\z)\to H_3(\SL_2(\z/p^r),\z)\to \RB(\z/p^r)\to 0
			\]
			A more complete result for finite local rings such as $\z/p^r$ was obtained by Mirzaii and Rojas in the (unpublished) work \cite[Theorem 6.2]{B-R--2025}.
		\end{enumerate}
	\end{exm}
	The next corollaries generalize Theorem 6.6 of \cite{B-E--2023} by removing the condition $H_2(B(A),\z)\simeq H_2(T(A),\z)$. 
	\begin{cor}
		Let $A$ be a ring satisfying the following conditions:
		\begin{enumerate}
			\item $\mu_2(A)=\{\pm 1\}$ and $-1\in\aa^2$;
			\item the complex $X_\bullet(A^2)\to\z$ is exact in dimension $<2$;
			\item $H_3(B(A),\z)\simeq H_3(T(A),\z)$.
		\end{enumerate}
		Then there exists an exact sequence
		\[
		H_3(\SM_2(A),\z)\to H_3(\SL_2(A),\z)\to \RB(A)\to 0
		\]
	\end{cor}
	\begin{proof}
		If $\mu_2(A)=\{\pm 1\}$ and $-1\in\aa^2$, then $\im(\gamma)=\im(\alpha)$ by \cite[Corollary 6.4]{B-E--2023}, and we may apply the second part of Theorem \ref{thm-main-1}.
	\end{proof}
	
	If $H_3(B(A),\z)\simeq H_3(T(A),\z)$, then $\coker(\alpha)\simeq H_3(\SL_2(A),\SM_2(A);\z)$; this follows from the analysis of the relative spectral sequence in \cite[\S 8]{B-E--2023}. The next corollary is a slight generalization of Theorem 8.2 of \cite{B-E--2023}, in which the condition $H_2(B(A),\z)\simeq H_2(T(A),\z)$ is no longer needed.
	\begin{cor}
		Let $A$ be a ring satisfying the following conditions:
		\begin{enumerate}
			\item the complex $X_{\bullet}(A^2)\to\z$ is exact in dimension $<2$;
			\item $H_3(B(A),\z)\simeq H_3(T(A),\z)$.
		\end{enumerate}
		Then there exists an exact sequence
		\[
		\II_A\otimes\mu_2(A)\to H_3(\SL_2(A),\SM_2(A);\z)\to \frac{\RP_1(A)}{\II_A\psi_1(-1)}\to 0
		\]
		If, in addition, $\im(\gamma)\subseteq\im(\alpha)$, then
		\[
		H_3(\SL_2(A),\SM_2(A);\z)\simeq\frac{\RP_1(A)}{\II_A\psi_1(-1)}
		\]
		and if $\im(\gamma)=\im(\alpha)$, then
		\[
		H_3(\SL_2(A),\SM_2(A);\z)\simeq\RP_1(A)
		\]
	\end{cor}
	\begin{proof}
		This is a direct consequence of Theorem \ref{thm-main} and of the observation made in the paragraph above.
	\end{proof}
	\begin{rem}
		In particular, if $A$ satisfies conditions (1) and (2) of the corollary above and, in addition, $\mu_2(A)=\{\pm 1\}$ and $-1\in\aa^2$, then $\im(\gamma)=\im(\alpha)$ and hence
		\[
		H_3(\SL_2(A),\SM_2(A);\z)\simeq \RP_1(A),
		\]
		which is precisely the second part of \cite[Theorem 8.2]{B-E--2023}. 
	\end{rem}
	As further examples, we obtain the isomorphisms
	\[
	H_3(\SL_2(\z\mth{2}),\SM_2(\z\mth{2});\z)\simeq \RP_1(\z\mth{2})
	\]
	and
	\[   
	H_3(\SL_2(\z\mth{10}),\SM_2(\z\mth{10});\z)\simeq \RP_1(\z\mth{10}).
	\]
	For $\F_q$ with $|\F_q|>81$ we have
	\[
	H_3(\SL_2(\F_q),\SM_2(\F_q);\z)\simeq \RP_1(\F_q).
	\]
	for $\z/p^r$ with $p>7$ we have,
	\[
	H_3(\SL_2(\z/p^r),\SM_2(\z/p^r);\z)\simeq \RP_1(\z/p^r).
	\]
	and
	\[
	H_3(\SL_2(\z),\SM_2(\z);\z)\simeq \RP_1(\z).
	\]
	\begin{rem}\label{lack-on-zm}
		The only obstruction to producing further examples of the form $A_m=\z\mth{m}$ seems to lie in the condition 
		\[
		H_3(B(A),\z)\simeq H_3(T(A),\z).    
		\]
		
		In Proposition \ref{prp-sm2-sl2}, this condition is used to guarantee that $E^\infty_{0,3}$ is a quotient of $H_3(T(A),\z)$. An alternative way of ensuring this is to require that $\im(d^2_{2,2})$ annihilate the summand $\SS_3$ of $H_3(B(A),\z)$.
		
		Along these lines, Proposition \ref{prp-hb-ht-am} (3) suggests further examples of rings $\z\mth{m}$ with $r=2$. For instance, the ring $\z\mth{290}$ satisfies all the conditions except the isomorphism above, and all that is needed is an element $x\in E^2_{2,2}$ such that $d^2_{2,2}(x)$ generates $\SS_3\simeq\z/3$. A similar situation occurs for the ring of Eisenstein integers $\z[\zeta]$. In fact, an analysis of the Lyndon/Hochschild-Serre spectral sequence of the extension
		\[
		1\to N(\z[\zeta]) \to B(\z[\zeta])\to T(\z[\zeta])\simeq\lan\zeta\ran\to 1
		\]
		shows that $\SS_3$ fits into the exact sequence
		\[
		0\to \z/6\to\SS_3\to \z/3\to 0.
		\]
	\end{rem}
	\section{A refined Bloch-Wigner sequence over local fields}
	The aim of this section is to obtain a refined Bloch-Wigner exact sequence for a non-archimedean local field $F$ with $\char(F)\neq 2$. We further assume that $F$ is non-dyadic, i.e., that $\char(\bar{F})\neq 2$, where $\bar{F}$ denotes the residue field. Let $\mathcal{O}_F$ be the ring of integers of $F$ and let $U=\mathcal{O}_F^\times$. It is well known that $\GG_F=\{1,\lan u\ran,\lan \pi\ran,\lan u\pi\ran\}$, where $\pi$ is a {\it uniformizer} and $u\in U$ is a unit whose reduction is not a square in $\bar{F}$. Since $\mu(F)$ is finite cyclic, let $\zeta$ be a generator and set $2N=|\mu(F)|$. We may take $u=\zeta$, or $u=-1$ if $-1$ is not a square in $F$, in which case $N$ is odd.
	
	Let $A$ be a domain with $\char(A)\neq 2$. Recall the exact sequence
	\[
	\begin{tikzcd}
		\II_A\otimes\mu_2(A)\ar[r,"\gamma"] & E^2_{2,1}\ar[r,"\delta"] & \RP_1(A) \ar[r] & 0
	\end{tikzcd}
	\]
	If $\II_A^2\otimes\mu_2(A)\subseteq\ker(\gamma)$ and $H_2(B(A),\z)\simeq H_2(T(A),\z)$ (i.e., $\SS_2=0$), then the differential $d^2_{2,2}$ yields the commutative diagram
	\[
	\begin{tikzcd}
		&\II_A\otimes\mu_2(A)\ar[r,"\gamma"]\ar[d] & E^2_{2,1}\ar[r,"\delta"]\ar[d,"d^2_{2,1}"] & \RP_1(A) \ar[r]\ar[d] & 0\\
		0 \ar[r] & \aa\wedge(-1) \ar[r]& \aa\wedge\aa \ar[r] & \frac{\aa\wedge\aa}{\aa\wedge(-1)} \ar[r] & 0
	\end{tikzcd}
	\]
	and the snake lemma yields the exact sequence
	\[
	\begin{tikzcd}
		\mathcal{K}\ar[r] & E^{\infty}_{2,1}\ar[r] & \RB(A)\ar[r] & 0.
	\end{tikzcd}
	\]
	
	Note that, by the proof of Corollary \ref{ker-gam}, every element of $\mathcal{K}$ is of the form $K+\Lan\nu\Ran\otimes(-1)$, where $K\in\II^2_A\otimes\mu_2(A)$ and $\nu$ is an $l$-th root of unity with $l$ even. If $\II^2_A\otimes\mu_2(A)\subset\ker(\gamma)$ and $\mu(F)$ is finite cyclic with generator $\zeta$, then we obtain the exact sequence
	\[
	\begin{tikzcd}
		\lan\Lan\zeta\Ran\otimes(-1)\ran\ar[r,"\gamma"] & E^{\infty}_{2,1}\ar[r,"\delta"] & \RB(A)\ar[r] & 0
	\end{tikzcd}
	\]
	Note that if $-1\notin\aa^2$, then $\lan\zeta\ran=\lan-1\ran$, and in the spectral sequence $\hat{E}^1_{p,q}\implies H_3(\SM_2(A),\z)$ the term $\hat{E}^{\infty}_{2,1}$ is generated by $\lan-1\ran$. We thus obtain the diagram
	\begin{equation}
		\begin{tikzcd}
			& \hat{E}^{\infty}_{2,1}\ar[d,"\alpha:=\inc_*"] & &\\
			\lan\Lan-1\Ran\otimes (-1)\ran\ar[r,"\gamma"]& E^{\infty}_{2,1} \ar[r,"\delta"] & \RB(A)\ar[r] & 0.
		\end{tikzcd}
	\end{equation}
	in which $\im(\gamma)=\im(\alpha)$. Applying the same reasoning as in Proposition \ref{prp-sm2-sl2}, we have proved the following result.
	\begin{prp}
		\label{prp-sm2-sl2-local}
		Let $A$ be a domain with $\char(A)\neq2$ such that
		\begin{enumerate}
			\item the complex $X_{\bullet}(A^2)\to\z$ is exact in dimension $<2$;
			\item $H_i(B(A),\z)\simeq H_i(T(A),\z)$ for $i=2,3$;
			\item $\mu(A)$ is finite cyclic and $-1\notin\aa^2$;
			\item $\II^2_{A}\otimes\mu_2(A)\subset\ker(\gamma)$.
		\end{enumerate}
		Then we have the exact sequence
		\[
		\begin{tikzcd}
			H_3(\SM_2(A),\z)\ar[r] & H_3(\SL_2(A),\z)\ar[r] & \RB(A) \ar[r] & 0.
		\end{tikzcd}
		\]
	\end{prp}
	This is the case for a non-dyadic local field $F$ with $\char(F)\neq 2$ and $-1\notin {F^{\times}}^2$ (the case in which $-1$ is a square follows from \cite[Corollary 6.4]{B-E--2023}). Indeed, conditions (1) and (2) hold because such fields are infinite, and condition (3) is a well-known fact about local fields. As for condition (4), take any product $\Lan a\Ran\Lan b\Ran\otimes(-1)$; the only relevant case is $a=-1$ and $b=\pi$, where $\pi$ is a uniformizer. In this case, Lemma \ref{alpha-gam-chi} (3) and (7) give
	\[
	\gamma(\Lan -1\Ran\Lan \pi\Ran\otimes(-1))=\Lan \pi\Ran\Lan -1\Ran\chi_1=0.
	\]
	Consequently, every such local field satisfies the above exact sequence.\\
	
	\begin{rem}
		\label{lf-classes}
		Note that in any non-dyadic local field we have $F^\times\wedge F^\times=\lan\pi\wedge\zeta\ran+2(F^\times\wedge F^\times)$, where $\zeta$ is a generator of $\mu(F)$; moreover, if $-1\notin {F^\times}^2$, we may replace $\zeta$ by $-1$. This is obtained simply by passing to the square classes of $a$ and $b$ in any element $a\wedge b$.
	\end{rem}
	
	We now analyze $H_3(\SM_2(F),\z)$ for a local field $F$ with $-1\notin {F^\times}^2$, using the spectral sequence $\hat{E}^1_{p,q}\implies H_3(\SM_2(F),\z)$. By \cite[Lemma 5.1]{B-E--2023}, $\displaystyle \hat{E}^1_{2,2}=H_2(\SM_2(F),\z)\simeq\frac{F^\times\wedge F^\times}{F^\times\wedge(-1)}$, and the map $\hat{d}^1_{2,2}$, which comes from the transfer map, satisfies $\overline{a\wedge b}\mapsto 2(a\wedge b)$. Lemma \ref{ker-wed-2} implies that $\hat{d}^1_{2,2}$ is injective, so that $\hat{E}^{\infty}_{2,2}=\hat{E}^2_{2,2}=0$. Since $\hat{d}^1_{1,2}=0$, the group $\hat{E}^{\infty}_{1,2}=\hat{E}^2_{1,2}$ is given by 
	\[
	\frac{\hat{E}^1_{1,2}}{\im(\hat{d}^1_{2,2})}=\frac{F^\times\wedge F^\times}{2(F^\times\wedge F^\times)}\simeq\lan\overline{\pi\wedge-1}\ran.
	\]
	
	In order to determine $\hat{E}^2_{0,3}$, we need some information about the group $H_3(F^\times,\z)$.
	
	Let $B$ be an abelian group and let $\sigma_1:\tors(B,B)\arr \tors(B,B)$
	be the map obtained by interchanging the two copies of $B$. It is not difficult to show that $\sigma_1$ is induced by
	the involution $B\otimes B \arr B\otimes B$, $a\otimes b\mapsto -b\otimes a$.
	
	Let $\Sigma_2'=\{1, \sigma'\}$ be the symmetric group of order $2$. Consider the following action of 
	$\Sigma_2'$ on $\tors(B,B)$:
	\[
	(\sigma', x)\mapsto -\sigma_1(x).
	\]
	
	\begin{prp}\label{H3B}
		For any abelian group $B$ we have the exact sequence
		\[
		\begin{array}{c}
			0 \arr \bigwedge_\z^3 B \arr H_3(B,\z) \arr \tors(B,B)^{\Sigma_2'} \arr 0,
		\end{array}
		\]
		where the homomorphism on the right is obtained from the composite
		\[
		H_3(B,\z) \overset{{\Delta_B}_\ast}{\larr } H_3(B\oplus B,\z) \arr \tors(B,B),
		\]
		$\Delta_B$ being the diagonal map $B \arr B\oplus B$, $b \mapsto (b,b)$.
	\end{prp}
	\begin{proof}
		See \cite[Lemma~5.5]{suslin1991}, \cite[Section 6]{breen1999}.
	\end{proof}
	
	When $B=F^\times$ we have $\tors(F^\times,F^\times)^{\Sigma'_2}=\tors(\mu(F),\mu(F))=\mu(F)$: the first equality holds because $F$ is a domain, and the second because $\mu(F)$ is finite. The map $\hat{d}^1_{1,3}$ induces the commutative diagram
	\[
	\begin{tikzcd}
		0 \ar[r] & \bigwedge^3F^\times \ar[r]\ar[d,"2"] & H_3(F^\times,\z)\ar[r]\ar[d,"\hat{d}^1_{1,3}"] & \mu(F)\ar[r]\ar[d,"0"] & 0\\
		0 \ar[r] & \bigwedge^3F^\times \ar[r] & H_3(F^\times,\z)\ar[r] & \mu(F)\ar[r] & 0
	\end{tikzcd}
	\]
	Indeed, $\hat{d}^1_{1,3}=w_*-\id_*$, where $w_*$ is induced by conjugation by the matrix $w$, and for any element $a\in F^\times$ this action sends $a\mapsto a^{-1}$; this implies that the left-hand map is multiplication by $2$. On the other hand, any generator $\zeta$ of the finite cyclic group $\mu(F)$ has a preimage $\theta_{\zeta}$ represented by the cycle (where $2N=|\mu(F)|$)
	\[
	\theta_{\zeta}=\sum^{2N-1}_{i=0}[\zeta|\zeta^i|\zeta]
	\]
	It is not difficult to show that $w\theta_{\zeta}-\theta_{\zeta}$ vanishes; in fact, it is the differential of
	\[
	\sum^{2N-1}_{i=0}\left\{[w|\zeta|\zeta^i|\zeta]-[\zeta^{-1}|w|\zeta^i|\zeta]+[\zeta^{-1}|\zeta^{-i}|w|\zeta]-[\zeta^{-1}|\zeta^{-i}|\zeta^{-1}|w]+[\zeta|\zeta^{-1}|\zeta^{-i}|\zeta^{-1}]-[\zeta|\zeta^i|\zeta|\zeta^{-1}]\right\}.
	\]
	
	This implies that the right-hand vertical map is $0$. The left-hand vertical map is surjective, since $\GG_F$ has only four elements: it suffices to pass to square classes in each factor of a generator and then use linearity. By the snake lemma we obtain $\hat{E}^2_{0,3}=\mu(F)$, and since $\hat{E}^2_{2,2}=0$, it follows that $\hat{E}^{\infty}_{0,3}=\mu(F)$. Moreover, $\hat{E}^2_{2,1}=\GG_F$ and $\hat{d}^2_{2,1}(\lan a\ran)=a\wedge(-1)$ (see \cite[\S 5]{B-E--2023}); hence, by Lemma \ref{ker-gam}, $\hat{E}^3_{2,1}=\ker(\hat{d}^2_{2,1})$ is generated by $\lan\zeta\ran$, where $\zeta$ is a generator of $\mu(F)$, and since $-1$ is not a square in $F$, we have $\lan\zeta\ran=\lan-1\ran$. 
	
	Consider now the filtration
	\[
	0\subseteq \hat{F}_0\subseteq \hat{F}_1\subseteq \hat{F}_2= \hat{F}_3=H_3(\SM_2(F),\z)
	\]
	This yields two exact sequences
	\begin{equation}
		\label{sq-sm2-1}
		0 \arr \hat{F}_0=\mu_2(F)\arr \hat{F}_1\arr \hat{E}^{\infty}_{1,2}\arr 0
	\end{equation}
	and
	\begin{equation}
		\label{sq-sm2-2}
		0 \arr \hat{F}_1\arr H_3(\SM_2(F),\z)\arr \hat{E}^{\infty}_{2,1}\arr 0
	\end{equation}
	
	In \eqref{sq-sm2-1}, a lift of $\overline{\pi\wedge\zeta}$ is the element
	\begin{align*}
		\Pi_{\zeta}=&\bigg\{[w|\zeta|\pi]-[w|\pi|\zeta]-[\zeta^{-1}|w|\pi]+[\pi^{-1}|w|\zeta]+[\zeta^{-1}|\pi^{-1}|w]-[\pi^{-1}|\zeta^{-1}|w]+[\zeta|\zeta^{-1}|\pi^{-1}]\\
		&-[\pi|\pi^{-1}|\zeta^{-1}]-[\pi|\zeta|\zeta^{-1}\pi^{-1}]+[\zeta|\pi|\zeta^{-1}\pi^{-1}]\bigg\}\otimes 1
	\end{align*}
	Its image in $H_3(\SM_2(F),\z)$ under the map in \eqref{sq-sm2-2} is a $2$-torsion element. This implies that $2\Pi_{\zeta}=0$, and hence that \eqref{sq-sm2-1} splits. More precisely,
	\begin{align*}
		2\Pi_{\zeta}-c(-1,\pi,\zeta)=&d_4\big(\{[w|w|\zeta|\pi]-[w|w|\pi|\zeta]-[w|\zeta^{-1}|w|\pi]+[w|\pi^{-1}|w|\zeta]+[\zeta|w|w|\pi]-[\pi|w|w|\zeta]\\
		&+[w|\zeta^{-1}|\pi^{-1}|w]-[w|\pi^{-1}|\zeta^{-1}|w]-[\zeta|w|\pi^{-1}|w]+[\pi|w|\zeta^{-1}|w]+[\zeta|\pi|w|w]\\
		&-[\pi|\zeta|w|w]-[\zeta^{-1}|w|\zeta^{-1}|\pi^{-1}]+[\pi^{-1}|w|\pi^{-1}|\zeta^{-1}]+[w|\zeta|\zeta^{-1}|\pi^{-1}]\\
		&-[w|\pi|\pi^{-1}|\zeta^{-1}]+[\zeta^{-1}|\zeta|w|\pi^{-1}]-[\pi^{-1}|\pi|w|\zeta^{-1}]-[\zeta^{-1}|\zeta|\pi|w]+[\pi^{-1}|\pi|\zeta|w]\\
		&+[\pi^{-1}|w|\zeta|\zeta^{-1}\pi^{-1}]-[\zeta^{-1}|w|\pi|\pi^{-1}\zeta^{-1}]-[\pi^{-1}|\zeta^{-1}|w|\zeta^{-1}\pi^{-1}]\\
		&+[\zeta^{-1}|\pi^{-1}|w|\pi^{-1}\zeta^{-1}]-[w|\pi|\zeta|\zeta^{-1}\pi^{-1}]+[w|\zeta|\pi|\pi^{-1}\zeta^{-1}]\\
		&-[\zeta^{-1}|\pi^{-1}|\zeta\pi|w]+[\pi^{-1}|\zeta^{-1}|\pi\zeta|w]-[\zeta^{-1}|\pi^{-1}|\pi|\zeta]+[\pi^{-1}|\zeta^{-1}|\zeta|\pi]\\
		&-[\zeta|\zeta^{-1}\pi^{-1}|\pi|\zeta]+[\pi|\pi^{-1}\zeta^{-1}|\zeta|\pi]-[\pi|\zeta|\zeta^{-1}\pi^{-1}|\pi\zeta]+[\zeta|\pi|\pi^{-1}\zeta^{-1}|\zeta\pi]\\
		&+[\zeta|\zeta^{-1}|\pi^{-1}|\pi]-[\pi|\pi^{-1}|\zeta^{-1}|\zeta]-[\pi|\pi^{-1}|\pi|\zeta]+[\zeta|\zeta^{-1}|\zeta|\pi]\}\otimes 1\big)
	\end{align*}
	where $c(-1,\pi,\zeta)$ is the image of $-1\wedge\pi\wedge\zeta\in\bigwedge^3F^\times$, and this last element vanishes because $-1=\zeta^{|\mu(F)|/2}$. Note that in all the computations above we may replace $\zeta$ by $-1$ when $-1\notin{F^\times}^2$; thus the element $\Pi_{-1}$ is $2$-torsion as well.\\ 
	\begin{prp}
		Let $F$ be a non-archimedean, non-dyadic local field with $\char(F)\neq 2$. If $-1\notin{F^\times}^2$, then
		\[
		H_3(\SM_2(F),\z)\simeq \frac{\z}{(4N)\z}\oplus\frac{\z}{2\z}. 
		\]
		where $2N=|\mu(F)|$.
		%\begin{enumerate}
		%    \item If $-1\notin{F^\times}^2$, then
		%    \[
		%        H_3(\SM_2(F),\z)\simeq \frac{\z}{(4N)\z}\oplus\frac{\z}{2\z}. 
		%    \]
		%    where $2N=|\mu(F)|$.
		%    \item If $-1\in{F^\times}^2$, then
		%    \[
		%        H_3(\SM_2(F),\z)\simeq \mu(F)\oplus\frac{\z}{2\z}\oplus\frac{\z}{2\z}. 
		%    \]
		%\end{enumerate}
	\end{prp}
	\begin{proof}
		When $-1$ is not a square, the exact sequence \eqref{sq-sm2-2} reads
		\[
		\begin{tikzcd}
			0 \ar[r]& \mu(F)\oplus\lan\overline{\pi\wedge(-1)}\ran\ar[r] & H_3(\SM_2(F),\z)\ar[r] & \hat{E}^{\infty}_{2,1}=\lan\lan -1\ran\ran\ar[r] & 0.
		\end{tikzcd}
		\]
		Taking a lift of $\lan -1\ran$ on the right-hand side, we obtain the cycle
		\[
		\Gamma = \{[w|w|-1]-[w|-1|w]+[w|-1|-1]+[-1|w|w]-[-1|w|-1]+[-1|-1|w]\}\otimes 1.
		\]
		It is not difficult to show that
		\begin{align*}
			2\Gamma-[-1|-1|-1]\otimes 1=&d\big(\{[w|w|w|-1]+[-1|-1|w|w]-[-1|w|-1|w]+[w|-1|-1|w]\\
			&+[w|w|-1|-w]-[w|-1|w|-w]+[-1|w|w|-w]\}\otimes 1\big).
		\end{align*}
		
		Now let $\Theta_\zeta$ denote the image of $\theta_{\zeta}$ in $H_3(\SM_2(F),\z)$. One checks that $\Theta_{\zeta}+\Gamma$ generates a cyclic subgroup of order $4N$ (recall that $N$ is odd), while $\Pi_{-1}$ generates a subgroup of order $2$; this implies that
		\[
		H_3(\SM_2(F),\z)\simeq\frac{\z}{(4N)\z}\oplus\frac{\z}{2\z}.
		\]
	\end{proof}
	
	Note that $\mu(F)^{\tilde{}}\simeq\frac{\z}{(4N)\z}$ is the unique non-trivial extension of $\mu(F)$ by $\z/2$. We now return to the spectral sequence $E^1_{p,q}\implies H_{p+q}(\SL_2(F),\z)$.
	\begin{lem}
		\label{d122-surj}
		If $-1\notin {F^{\times}}^2$, then $d^1_{2,2}$ is surjective.
	\end{lem}
	\begin{proof}
		For such fields we know that $\GG_F=\{1,\lan u\ran,\lan\pi\ran,\lan u\pi\ran\}$, where $u$ is a non-square in the residue field.
		If $-1\notin {F^\times}^2$, we may take $u=-1$, and $\pi+1=r^2$ for some $r$ (by Hensel's lemma). Consider now the long exact sequence \eqref{long-eseq} and the element $[-1]\otimes\partial_3((\pmb{\infty},\pmb{0},\pmb{1},\pmb{\pi+1}))$, which represents an element of $H_1(\SL_2(F),Z_2)$. This element maps to $\Lan\pi+1\Ran\Lan-\pi\Ran\otimes(-1)\in\RR_F\otimes\mu_2(F)$, which is zero. Hence there exists an element $l(\pi)\in H_2(\SL_2(F),Z_1)$ which is represented in the bar resolution by
		\[
		l(\pi)=R_{\pi}\otimes\partial_2(X_{-\pi^{-1}})+H_{\pi}\otimes\partial_2(X_1)
		\]
		where
		\[
		R_{\pi}=[g_1^{-1}|-1]-[-1|g_1^{-1}]-[h_1^{-1}r|-1]+[-1|h_1^{-1}r],\quad H_{\pi}=[r|-1]-[-1|r].
		\]
		where $h_{a}=\mtxx{1}{a^{-1}}{0}{1}$ and $g_a=\mtxx{0}{1}{-1}{a}$ for $a\in F^{\times}$. We shall show that $d^1_{2,2}(l(\pi))=-1\wedge\pi\in H_2(T(F),\z)\simeq H_2(\SL_2(F),X_1)$ in the spectral sequence $E^1_{p,q}$. This image is represented by
		\[
		(g_{-\pi^{-1}}^{-1}-h_{-\pi}^{-1}+1)R_{\pi}+(g_1^{-1}-h_1^{-1}+1)H_{\pi})\otimes(\pmb{\infty},\pmb{0})
		\]
		This is an element of $H_2(\SL_2(F),X_1)\simeq H_2(T(F),\z)$. We shall use the section trick; to this end, we pass the element above to the standard resolution: 
		\begin{align*}
			(g_{-\pi^{-1}}^{-1}-h_{-\pi}^{-1}+1)&((-1,g_1^{-1},-g_1^{-1})-(1,-1,-g_1^{-1})-(1,h_1^{-1}r,-h_1^{-1}r)+(1,-1,-h_1^{-1}r)))\otimes(\pmb{\infty},\pmb{0})\\
			+(g_1^{-1}-h_1^{-1}+1)&((1,r,-r)-(1,-1,-r))\otimes(\pmb{\infty},\pmb{0})
		\end{align*}
		Take a section $s:T(F)\backslash\SL_2(F)\to \SL_2(F)$ of the canonical projection $p:\SL_2(F)\to T(F)\backslash\SL_2(F)$, given by
		\[
		s(T(F)\mtxx{a}{b}{c}{d}):=\begin{cases}
			\mtxx{1}{ba^{-1}}{ca}{da}, & a\neq 0,\\
			\mtxx{0}{1}{-1}{bd}, & a = 0
		\end{cases}
		\]
		We define a map $g\mapsto\overline{g}:=g(s\circ p(g))^{-1}$, which allows us to pass from the standard $\SL_2(F)$-resolution to a standard $T(F)$-resolution. Applying this map and simplifying, we obtain
		\begin{align*}
			((-\pi^{-1},-(\pi^{-1}+1),\pi^{-1}+1)&-(-\pi^{-1},\pi^{-1},\pi^{-1}+1)\\
			-(1,\pi^{-1}+1,-(\pi^{-1}+1))&+(1,-1,-(\pi^{-1}+1)))\otimes(\pmb{\infty},\pmb{0})
		\end{align*}
		which, on passing to the bar resolution, becomes $([1+\pi|-1]-[-1|1+\pi]-[\pi^{-1}+1|-1]+[-1|\pi^{-1}+1])\otimes(\pmb{\infty},\pmb{0})$; this represents the element $(1+\pi)\wedge(-1)-(\pi^{-1}+1)\wedge(-1)=\pi\wedge(-1)$ in $H_2(\SL_2(F),X_1)\simeq H_2(T(F),\z)\simeq\bigwedge^2F^{\times}$.\\
		
		Finally, Remark \ref{lf-classes} shows that $a\wedge b=n\{\pi\wedge(-1)\}+2K$ for any $a\wedge b\in\bigwedge^2 F^{\times}$, with $K\in\bigwedge^2F^{\times}$, and it is straightforward to check that the element $([u|v]-[v|u])\otimes\{(\pmb{\infty},\pmb{0})+(\pmb{0},\pmb{\infty})\}$, which represents an element of $E^1_{2,2}$, maps to $2(u\wedge v)$. This completes the proof.
	\end{proof} 
	\begin{rem}
		\label{rem-d211-gen}
		More generally, if $A$ is a universal $\GE_2$-ring consider the long exact sequence \eqref{long-eseq}. For $a\in\WW_A$, the cycle
		\[
		p_{-1}^+[-1]\otimes\partial_3(\pmb{\infty},\pmb{0},\pmb{1},\pmb{a})+\Lan1-a\Ran[-1]\otimes\partial_3((\pmb{\infty},\pmb{0},\pmb{1},\pmb{a})+(\pmb{0},\pmb{\infty},\pmb{1},\pmb{a}))
		\]
		where $p_{-1}^+=\lan -1\ran+1$, represents an element $\tilde{g}(a)$ which maps to zero in $H_1(\SL_2(A),X_2)\simeq\RR_A\otimes\mu_2(A)$. Hence there exists an element
		\begin{align*}
			g(a)=&([g_1^{-1}|-1]-[-1|g_1^{-1}]-[1-a|-1]+[-1|1-a])\otimes X_{(1-a)^{-1}}\\
			&-([h_1^{-1}|-1]-[-1|h_1^{-1}]-[1-a|-1]+[-1|1-a])\otimes X_{a(1-a^{-1})}\\
			&+([g_{-1}^{-1}|-1]-[-1|g_{-1}^{-1}]-[w|-1]+[-1|w])\otimes X_{-(1-a)^{-1}}\\
			&-([h_{-1}^{-1}|-1]-[-1|h_{-1}^{-1}]-[wa|-1]+[-1|wa])\otimes X_{-a(1-a)^{-1}}\\
			&-([wa|-1]-[-1|wa])\otimes X_{-a}+([w|-1]-[-1|w])\otimes X_{-1}
		\end{align*}
		in $H_2(\SL_2(A),Z_1)$. The same reasoning, together with the section trick used in the proof above (with the same map $s$), shows that $d^1_{2,2}(g(a))=-1\wedge (1-a)$. Setting $a=1-\pi$ recovers the result above, but the argument applies to other universal $\GE_2$-rings with $|\GG_A|=4$ in which $-1$ is not a square, for example $\z\half$.  
	\end{rem}
	The lemma above implies that $E^{\infty}_{1,2}=E^2_{1,2}=0$. The spectral sequence yields a filtration
	\[
	0\subseteq F_0 = F_1\subseteq F_2 = F_3=H_3(\SL_2(F),\z)
	\]
	and comparison with the spectral sequence for $\SM_2(F)$ gives the commutative diagram
	\[
	\begin{tikzcd}
		0\ar[r] & \mu(F)\ar[r]\ar[d] & \hat{F}_1 \ar[r]\ar[d] & \hat{E}^{\infty}_{1,2}=\lan\overline{\pi\wedge(-1)}\ran\ar[r]\ar[d] & 0\\
		0\ar[r] & \frac{\mu(F)}{\im(d^2_{2,2})}\ar[r] & F_1 \ar[r] & E^{\infty}_{1,2}=0\ar[r,equal] & 0
	\end{tikzcd}.
	\]
	
	The upper row is split by $s(\overline{\pi\wedge(-1)})=\Pi_{-1}$, which yields the following commutative diagram. 
	\[
	\begin{tikzcd}
		\lan\overline{\pi\wedge(-1)}\ran\ar[r]\ar[d] & H_3(\SM_2(F),\z)\ar[d]\\
		0\ar[r] & H_3(\SL_2(F),\z)
	\end{tikzcd}
	\]
	Hence the image of $\Pi_{-1}$ under $H_3(\SM_2(F),\z)\to H_3(\SL_2(F),\z)$ is trivial, and the exact sequence \eqref{sm2-sl2-seq} (over the field $F$) yields the exact sequence
	\[
	\begin{tikzcd}
		\mu(F)^{\tilde{}}\ar[r] & H_3(\SL_2(F),\z)\ar[r] & \RB(F)\ar[r] & 0.
	\end{tikzcd}
	\]
	
	Finally, from the classical Bloch-Wigner exact sequence over the algebraic closure $\pmb{F}$ of $F$, we have the commutative diagram
	\[
	\begin{tikzcd}
		&\mu(F)^{\tilde{}}\ar[r]\ar[d,hook] & H_3(\SL_2(F),\z)\ar[r]\ar[d] & \RB(F)\ar[r]\ar[d] & 0\\
		0 \ar[r] & \mu(\pmb{F})^{\tilde{}}\ar[r] & K_3(\pmb{F})\simeq H_3(\SL_2(\pmb{F}),\z)\ar[r] & \BB(\pmb{F})\ar[r] & 0
	\end{tikzcd}.
	\]
	
	The injectivity of the left-hand vertical map implies the injectivity of the left-hand map in the upper row, and we finally obtain the following theorem.
	\begin{thm}
		Let $F$ be a non-dyadic local field with $\char(F)\neq 2$ and $-1\notin {F^\times}^2$. There is a refined Bloch-Wigner exact sequence
		\[
		0\to\mu(F)^{\tilde{}}\to H_3(\SL_2(F),\z)\to\RB(F)\to 0.
		\]
		where $\mu(F)^{\tilde{}}$ is the unique non-trivial extension of $\mu(F)$ by $\z/2$.
	\end{thm}
	\begin{rem}
		\label{rem-z12}
		The reasoning above applies to any ring with $\GG_A=\{1,\lan-1\ran,\lan p\ran,\lan-p\ran\}$ (where $1-p$ is a unit), for example $\z\half$; an analysis of the spectral sequence for $\SM_2$ then shows that
		\[
		H_3(\SM_2(\z\half),\z)\simeq\frac{\z}{4\z}\oplus\frac{\z}{2\z}.
		\]
		
		The surjectivity of $d^1_{2,2}$ is established in Remark \ref{rem-d211-gen}. Using the same reasoning for the injectivity of the left-hand map (this time passing to $\q$), we obtain the refined Bloch-Wigner exact sequence
		\[
		0\to\frac{\z}{4\z}\to H_3(\SL_2(\z\half),\z)\to\RB(\z\half)\to0
		\]
		which confirms the results of Coronado and Hutchinson in \cite[\S 8]{C-H2022}.
	\end{rem}
	%%%%%%%%%%%%%%%%%%%%%%%%%%%%%%%%%%%%%%%%%%%%%%%%%%%%%%%%
	\begin{thebibliography}{99}
		%%%%%%%%%%%%%%%%%%%%%%%%%%%%%%%%%%%%%%%%%%%%%%%%%%%%%%%%
		
		%\bibitem{adem-naffah1998}
		%Adem, A., Naffah, N. On the cohomology of $\SL_2(\z[1/p])$. In Geometry and cohomology in group 
		%theory (Durham, 1994), volume 252 of London Math. Soc. Lecture Note Ser., pages 1--9. Cambridge Univ. P
		%ress, Cambridge, 1998
		
		\bibitem{breen1999}
		Breen, L. On the functorial homology of abelian groups. 
		Journal of Pure and Applied Algebra {\pmb 142} (1999) 199--237.
		%\bibitem{bass-tate1973} 
		%Bass, H., Tate, J. The Milnor ring of a global field. Lecture Notes in Math., Vol. {\bf 342},  349--446 (1973)
		
		\bibitem{bloch2000}
		Bloch, S. J. Higher Regulators, Algebraic $K$-Theory and Zeta Functions of Elliptic Curves. 
		CRM Monograph Series. 11. Providence, RI: American Mathematical Society (AMS), 2000.
		
		\bibitem{brown1994}
		Brown, K. S.: Cohomology of Groups. Graduate Texts in Mathematics, 87. Springer-Verlag, New York (1994)
		
		\bibitem{cohn1966}
		Cohn, P. M. On the structure of the $\GL_2$ of a ring. Inst. Hautes \'Etudes Sci. Publ. Math. {\bf 30} (1966), pp. 5--33
		
		\bibitem{cohn1968}
		Cohn P. M. A presentation of $\SL_2$ for Euclidean imaginary quadratic number fields. Mathematika {\bf 15(2)}. (1968), pp. 156--163.
		
		\bibitem{C-H2022}
		Coronado, R. C.,  Hutchinson, K. Bloch Groups of Rings. https://arxiv.org/abs/2201.04996
		
		\bibitem{dupont-sah1982}
		Dupont, J. L., Sah, C. Scissors congruences. II. J. Pure Appl. Algebra  {\bf 25}  (1982), no. 2, 159--195.
		
		%\bibitem{gsz2016}
		%Gille S., Scully S., Zhong, C. Milnor-Witt $K$-groups of local rings. Adv. Math. {\bf 286}  (2016) , 729--753
		
		%\bibitem{HM1973}
		%Husemoller, D., Milnor, J. Symmetric bilinear forms. Springer-Verlag, Berlin, Heidelberg, New York, 1973
		
		\bibitem{hutchinson-2013}
		Hutchinson, K. A Bloch-Wigner complex for $\SL_2$. J.~$K$-Theory {\bf 12}, no. 1, 15--68 (2013)
		
		%\bibitem{hutchinson2013}
		%Hutchinson, K. A refined Bloch group and the third homology of $\SL_2$ of a field. J. Pure Appl. Algebra {\bf 217}, 2003--2035 (2013)
		
		%\bibitem{hutchinson2015}
		%Hutchinson, K. On the low-dimensional homology of $\SL_2(k[t, t^{-1}])$. J. Algebra {\bf 425}, 324--366 (2015)
		
		\bibitem{hutchinson2017}
		Hutchinson, K. The third homology of $\SL_2$ of local rings. Journal of Homotopy and Related Structures {\bf 12}, 931--970 (2017)
		
		\bibitem{hutchinson2022}
		Hutchinson, K. $\GE_2$-rings and a graph of unimodular rows. J. Pure Appl. Algebra {\bf 226}, (2022), no. 10, 107074.
		
		\bibitem{knudson2001}
		Knudson, Kevin P. Homology of Linear Groups. Birk\"hauser Basel, 2001.
		%\bibitem{hutchinson-2017}
		%Hutchinson, K. The third homology of $\SL_2$ of fields with discrete valuation. J. Pure Appl. Algebra
		%{\bf 221} (5), 1076--1111 (2017)
		
		%\bibitem{hutchinson2019}
		%Hutchinson, K.: The third homology of $\SL_2(\q)$. Preprint. arxiv.org/abs/1906.11650
		
		
		
		%\bibitem{hutchinson-tao2009}
		%Hutchinson, K., Tao, L. The third homology of the special linear group of a field. J. Pure Appl.
		%Algebra {\bf 213} (2009), no. 9, 1665--1680
		
		%\bibitem{hutchinson-tao2010}
		%Hutchinson, K., Tao, L. Homology stability for the special linear group of a field and Milnor-Witt $K$-theory. Documenta Mathematica (2010) Volume: Extra Vol., 267--315
		
		%\bibitem{hmm2022}
		%Hutchinson, K., Mirzaii, B.,  Mokari, F. Y. The homology of $\SL_2$ of discrete valuation rings. Advances in Mathematics {\bf 402} (2022) 108313
		
		%\bibitem{lam2005}
		%Lam, T. Y. Introduction to Quadratic Forms over Fields. Graduate Studies in Mathematics. Vol. 67. American Mathematical Society 2005
		
		
		%\bibitem{lichenbaum1989}
		%Lichtenbaum, S. Groups related to scissors-congruence groups.Algebraic $K$-theory and algebraic 
		%number theory, 151--157, Contemp. Math., 83, Amer. Math. Soc., Providence, RI, 1989.
		
		%\bibitem{loday1976} 
		%Loday, J. L. $K$-th\'eorie alg\'ebrique et repr\'esentations de groupes. Ann. Sci. 
		%\'Ecole Norm. Sup. (4) {\bf 9} (1976), no. 3, 309--377.
		
		%\bibitem{maclane1963}  Mac Lane, S. Homology. New York; Springer-Verlag, Berlin-G\"ottingen-Heidelberg 1963.
		
		%\bibitem{maz2005}
		%Mazzoleni A. A new proof of a theorem of Suslin. $K$-Theory {\bf 35} (3-4) 199--211 (2005)
		
		%\bibitem{merkurjev-suslin1990}
		%Merkurjev, A. S., Suslin, A. A.: The group $K_3$ for a field.
		%(Russian) Izv. Akad. Nauk SSSR Ser. Mat. 54 (1990), no. 3, 522--545; translation in
		%Math. USSR-Izv. {\bf 36} , no. 3, 541--565 (1991)
		
		%\bibitem{mirzaii-2008}
		%Mirzaii, B. Third homology of general linear groups. J. Algebra {\bf 320} (2008), no. 5, 1851--1877.
		
		%\bibitem{mirzaii2011}
		%Mirzaii, B. Bloch-Wigner theorem over rings with many units. Math. Z. {\bf 268}, 329--346 (2011). Erratum to: Bloch-Wigner theorem over rings with many units. Math. Z. {\bf 275}, 653--655 (2013)
		
		%\bibitem{mm2015}
		%Mirzaii, B.,  Mokari, F. Y. A Bloch-Wigner theorem over rings with many units II.
		%Journal of Pure and Applied Algebra {\bf 219}, 5078--5096 (2015)
		
		%\bibitem{mirzaii2017}
		%Mirzaii, B. A Bloch-Wigner exact sequence over local rings. Journal of Algebra {\bf 476}, 459--493 (2017)
		
		%\bibitem{B-E-2023}
		%Mirzaii, B., Torres Pérez, E. A refined Bloch-Wigner exact sequence in characteristic 2. 
		%Journal of Algebra {\bf 657} (2024), no. 1, 141--158.
		
		
		\bibitem{B-E--2023}
		Mirzaii, B., Torres Pérez, E.
		A refined scissors congruence group and the third homology of $\SL_2$. Journal of Pure and Applied Algebra, {\bf 228} (2024), no. 6, 107615, 1--28.
		
		\bibitem{B-E-2025}
		Mirzaii, B., Torres Pérez, E.
		The third homology of projective special linear group of degree two. Journal of Pure and Applied Algebra, {\bf 229} (2025), no. 6, 107965.
		
		\bibitem{B-E--2026}
		Mirzaii, B. Torres Pérez, E. The low-dimensional homology groups of elementary group of degree two. Journal of Homotopy and Related Structures (2026).
		
		%\bibitem{B-E--2025}
		%Mirzaii, B. Torres Pérez, E. The low-dimensional homology of projective linear group of degree two. Doc. Math. 30 (2025), no. 5, pp. 1157–1199.
		
		%\bibitem{nes-suslin1990}
		%Nesterenko Yu. P., Suslin A. A. Homology of the general linear group over a local ring and Milnor's $K$-theory. Math. USSR-Izv. {\bf 34} (1990), no. 1, 121--145.
		
		%\bibitem{sah1989}
		%Sah, C. Homology of classical Lie groups made discrete. III. J. Pure Appl. Algebra {\bf 56} (1989), no. 3, 269--312.
		
		%\bibitem{jmorita1990}
		%On the group structure of rank one $K_2$ of some $\z_S$. Bull. Soc. Math. Belg., S\'er. A, {\bf 42} (1990), no. 3, pp. 561--575.
		
		%\bibitem{serre1980}
		%Serre, J.-P. Trees, Springer-Verlag, Berlin  (1980) 
		
		%\bibitem{sherman1989}
		%Sherman, C. K-theory of discrete valuation rings. ?Journal of Pure and Applied 
		%Algebra {\bf 61}  79--98, (1989)
		
		%\bibitem{srinivas1996}
		%Srinivas, V. Algebraic $K$-Theory. Second edition. Progress in Mathematics, 90. Birkh\"auser Boston, 1996.
		
		%\bibitem{suslin1985} 
		%Suslin, A. A. Homology of ${\rm GL}\sb{n}$, characteristic classes and Milnor $K$-theory. Proc. 
		%Steklov Math. {\bf 3} (1985), 207--225.
		\bibitem{B-R--2025}
		Mirzaii, B. Rojas. A. Schur multiplier of $\mathrm{SL}_2$ over finite commutative rings. https://arxiv.org/abs/2510.03946
		
		\bibitem{suslin1991}
		Suslin, A. A. $K\sb 3$ of a field and the Bloch group. Proc. Steklov Inst. Math. {\bf 183}, no. 4, 217--239 (1991)
		
		%\bibitem{vdkallen1977}
		%Van der Kallen, W. The $K\sb{2}$ of rings with many units. Ann. Sci. \'Ecole Norm. Sup. (4) 
		%{\bf 10}  (1977), no. 4, 473--515
		
		%\bibitem{vas1972}
		%Vaserstein L. N. On the group $\SL_2$ over Dedekind rings of arithmetic type. 
		%Math. USSR Sbornik {\bf 18} (1972), no. 2, 321--332
		
		%\bibitem{weibel1994}
		%Weibel, C. A. An Introduction to Homological Algebra. Cambridge Studies in Advanced Mathematics, 
		%38. Cambridge University Press, Cambridge, (1994).
		
		%\bibitem{weibel2013}
		%Weibel, C. A. The $K$-book: An introduction to algebraic $K$-theory. Graduate Studies in 
		%Mathematics, vol. 145. American Mathematical Society, Providence (2013)
		
		%\bibitem{wendt2018}
		%Wendt, M. Homology of $\SL_2$ over function fields I: parabolic subcomplexes. J. Reine Angew. Math. {\bf 739} (2018), 159--205
		
		%\bibitem{zikert2015}
		%Zickert, C. K. The extended Bloch group and algebraic K-theory. J. Reine Angew. Math. 
		%{\bf 704}, 21--54, (2015)
		
	\end{thebibliography}
	
\end{document}